\documentclass[11pt]{article}
\usepackage[utf8]{inputenc}
\usepackage[T1]{fontenc}
\usepackage{lmodern}
\usepackage{amsmath, amssymb, amsthm}
\usepackage{array}
\usepackage[margin=1in]{geometry}
\usepackage{microtype}
\PassOptionsToPackage{hyphens}{url}
\usepackage{xcolor}
\usepackage{hyperref}
\hypersetup{colorlinks=true, linkcolor=blue!50!black, urlcolor=blue!50!black, citecolor=blue!50!black,
  pdftitle={Differential transcendence of the Hahn-Exton q-cosine and the arithmetic of its zeros},
  pdfauthor={Vico Bonfioli}}

\theoremstyle{plain}
\newtheorem{theorem}{Theorem}
\newtheorem{lemma}[theorem]{Lemma}
\newtheorem{corollary}[theorem]{Corollary}
\newtheorem{proposition}[theorem]{Proposition}
\newtheorem{conjecture}[theorem]{Conjecture}
\theoremstyle{definition}
\newtheorem{definition}[theorem]{Definition}
\theoremstyle{remark}
\newtheorem{remark}[theorem]{Remark}

\newcommand{\Z}{\mathbb{Z}}
\newcommand{\Q}{\mathbb{Q}}
\newcommand{\C}{\mathbb{C}}
\DeclareMathOperator{\ord}{ord}
\DeclareMathOperator{\Cas}{Cas}

\title{Differential transcendence of the Hahn--Exton $q$-cosine\\ and the arithmetic of its zeros}
\author{Vico Bonfioli\thanks{Independent researcher. \texttt{vicobonfioli@gmail.com}. Code and
verification data: \url{https://github.com/ElVec1o/pythagorean-reflection-sequence}.}}
\date{\today}

\begin{document}
\maketitle

\begin{abstract}
Let $G(q,z)={}_1\phi_1(0;q;q^2,z)=\sum_{k\ge0}(-1)^kq^{k(k-1)}z^k/(q;q)_{2k}$, the Hahn--Exton (third
Jackson) $q$-cosine. We prove that for every $0<q<1$ the difference Galois group over $\C(z)$ of the
second-order $q$-difference equation satisfied by $G$ contains $\mathrm{SL}_2$, and that the equation
admits no $\sigma\delta$-integrability relation; via the parametrized difference Galois theory of
Hardouin--Singer \cite{HS} and Dreyfus--Hardouin--Roques \cite{DHR2018} these give a second route to the differential
transcendence of $G(q,\cdot)$ over $\C(z)$. That transcendence is itself known, and in stronger form:
it is a corollary of Adamczewski--Dreyfus--Hardouin \cite{ADH2021}, and the unsolvability of the
associated Riccati equation is Nishioka's \cite{Nishioka2016}; what is offered here is the explicit
Galois-theoretic computation, valid also at algebraic $q$. We then prove the connection constant
$C_1=\lim_{R\to\infty}G(q,-R)/\theta_{q^2}(R)=1/(q;q)_\infty$ against the divergent theta solution,
and deduce from it the regularized product of the positive zeros,
$\prod_{k\ge1}z_k(q)\,q^{2(k-1)}=(q;q^2)_\infty$. The product is a corollary of that constant;
an earlier attribution of the constant to Morita's connection formula \cite{Morita2011} is
withdrawn in \S\ref{sec:zeroproduct}, where the reasons are given. We then pose the arithmetic
question that these results frame: the transcendence of the zeros $z_k(q)$ at algebraic $q$. That
question specializes to the transcendence of a specific reflection-group growth constant $\beta_2$
(OEIS A396406), for which we prove an unconditional effective non-rationality bound.
\end{abstract}

\medskip
\noindent\textbf{Keywords.} Hahn--Exton $q$-cosine; third Jackson $q$-Bessel function; differential
transcendence; parametrized difference Galois theory; $q$-difference equation; zeros of
$q$-special functions; regularized product; partition integrality; effective irrationality.

\smallskip
\noindent\textbf{2020 Mathematics Subject Classification.} Primary 11J81, 33D15; Secondary 12H10,
33D45, 39A13, 11P84.

\section{Introduction}

The three Jackson $q$-analogues of the Bessel function are among the most-studied special functions of
$q$-series. The third goes back to Hahn \cite{Hahn1949} and, in the form now standard, to
Exton \cite{Exton1978}; the attribution and the modern normalization are those of
Koelink--Swarttouw \cite{KS} and Koornwinder--Swarttouw \cite{KS1992}. It has the $q$-cosine and
$q$-sine specializations
\[
   G(q,z):={}_1\phi_1(0;q;q^2,z)=\sum_{k\ge0}\frac{(-1)^kq^{k(k-1)}}{(q;q)_{2k}}z^k,
\]
which is, up to an explicit prefactor, $J^{(3)}_{-1/2}(\sqrt z/q;q^2)$; it is entire of order $0$ in
$z$, and its positive zeros $0<z_1(q)<z_2(q)<\cdots$ are real and simple \cite{KS}. The analytic theory
of these zeros (orthogonality, interlacing, asymptotics) is classical. Their \emph{arithmetic}
is, to our knowledge, untouched: nothing is known about the transcendence of $z_k(q)$ at algebraic
$q$, the $q$-analogue of Siegel's theorem that the zeros of the Bessel function $J_\nu$ are
transcendental.

This paper establishes the function-theoretic groundwork for that arithmetic and states the resulting
transcendence problem precisely. Our results are of three kinds.

\emph{Galois theory (\S\ref{sec:galois}).} We compute enough of the difference Galois structure of the
defining equation to conclude differential transcendence.
\begin{theorem}[known; Adamczewski--Dreyfus--Hardouin]\label{thm:main-transc}
For every $0<q<1$, the Hahn--Exton $q$-cosine $G(q,\cdot)$ satisfies no algebraic differential
equation over $\C(z)$.
\end{theorem}

\noindent This is \emph{not} new, and the known result is stronger. Adamczewski, Dreyfus and
Hardouin \cite[Thm.~1.2, case $Q$]{ADH2021} prove that a Puiseux-series solution of a linear
$q$-difference equation over $K=\bigcup_j\C(x^{1/j})$, with $q$ not a root of unity, is either in
$K$ or hypertranscendental over $K$; $G(q,\cdot)$ is entire and not in $K$, so the conclusion
follows with no hypothesis to verify, over a larger base field, and for all $q$ that are not roots
of unity. It is recorded here only because the route below is different and yields the explicit
Galois-theoretic statements of \S3, which \cite{ADH2021} does not compute. The unsolvability of the
associated Riccati equation for $J^{(3)}_\nu$, and the non-solvability of the Galois group, were
obtained for all $\nu$ by Nishioka \cite{Nishioka2016} for $q$ transcendental over $\Q$; the only
increment here is that algebraic $q$ are included.
The engine is an unconditional proof that the difference Galois group contains $\mathrm{SL}_2$
(Theorem~\ref{thm:sl2}), obtained by ruling out rank-$1$ submodules of every symmetric power through a
growth estimate on asymptotic-expansion coefficients, together with the absence of a
$\sigma\delta$-integrability relation (Proposition~\ref{prop:nonintegrable}).

\emph{A modular zero product (\S\ref{sec:zeroproduct}).} Writing $Q=q^2$,
\begin{theorem}\label{thm:zeroproduct-intro}
$\displaystyle\prod_{k\ge1}z_k(q)\,Q^{\,k-1}=(q;q^2)_\infty$, and $G(q,z)/\theta_Q(-z)\to1/(q;q)_\infty$
along the negative axis, where $\theta_Q(-z)=\sum_n(-1)^nq^{n(n-1)}z^n$.
\end{theorem}
The regularized product of the zeros is thus an explicit modular value: transcendental at algebraic
$q$ by Nesterenko's theorem \cite{Nesterenko}.

\emph{Arithmetic of the zeros (\S\ref{sec:arith}).} These results frame a clean conjecture.
\begin{conjecture}[Z]\label{conj:Z}
For algebraic $q\in(0,1)$, every positive zero $z_k(q)$ of $G(q,\cdot)$ is transcendental.
\end{conjecture}

\noindent Two remarks fix the shape of this conjecture. First, no blanket statement of this kind can
hold: the full Jacobi theta $\theta(z;q)=(z;q)_\infty(q/z;q)_\infty$ has zeros exactly at $z=q^{-k}$,
algebraic whenever $q$ is, so any hypothesis must exclude functions with a product formula, and it is
the failure of the triple product for $G$ that carries the conjecture. This is the structural
analogue of Siegel's exclusion of $2\lambda$ odd. Second, the classical model needs stating with
care, because the present object sits in the excluded case. Siegel \cite{Siegel1929} proves
that for rational $\lambda$ with $2\lambda$ \emph{not} an odd integer, every nonzero zero of
$J_\lambda$ is transcendental; here $2\lambda=-1$, which is exactly what that hypothesis removes.
For $\lambda=-\tfrac12$ one has $J_{-1/2}(x)\propto x^{-1/2}\cos x$, whose zeros are the odd
multiples of $\pi/2$, and their transcendence is Lindemann's theorem \cite{Lindemann1882}, not
Siegel's. So the classical analogue of the conjecture below is Lindemann's, and the $q$-deformation
is what makes it hard again (see also Lorch--Muldoon \cite{LorchMuldoon}). We have found no result, and no conjecture, in the
literature on the arithmetic nature of zeros of $q$-series; the $q$-Bessel zero literature is
uniformly classified as special functions, not number theory.
We record where Conjecture~\ref{conj:Z} stands relative to the $q$-analogue of the
Siegel--Shidlovskii method. The position is weaker than a margin: as \S\ref{sec:arith} sets out,
the theorem of \cite{AMV} yields linear independence over a fixed field containing $q$ and not
transcendence over $\overline\Q$, so no choice of parameters would make it apply to
Conjecture~\ref{conj:Z}; what the double-Pochhammer denominator $(q;q)_{2k}$ controls is the
bookkeeping ratio of \S\ref{sec:arith}, which is a factor $2$ past a boundary that is itself
already out of reach (Remark~\ref{rem:thetabarrier}). We also record an unconditional partial
result for the specific zero of number-theoretic interest.

The motivating instance is the growth constant $\beta_2=1.4916177871\ldots$ of the geodesic growth
series of a planar right-triangle reflection group (OEIS A396406), which by \cite{Bonfioli2}
equals $1/\sqrt{q^\ast}$ where
$q^\ast$ is the least positive root of $S(q):=G(q,2q(1-q))$; equivalently $z^\ast:=2q^\ast(1-q^\ast)=
z_1(q^\ast)$ is the smallest positive zero of $G(q^\ast,\cdot)$. Conjecture~\ref{conj:Z} implies
$\beta_2$ transcendental. We prove unconditionally (Proposition~\ref{prop:effective}) that $q^\ast$ has
no rational representation of denominator $\le10^{1044}$.

Every numerical claim below is regenerated by a script committed in the repository above, named at
the point of use, and each is reported with the agreement actually attained rather than with the
working precision. The scripts are \texttt{qcos\_identities.py} (symbolic identities),
\texttt{qcos\_lattice.py} (the integer zero lattice, exact in $\Z[[q]]$),
\texttt{qcos\_analytic.py} (the high-precision certificates, each run at two working precisions),
\texttt{qcos\_ledger.py} (the Diophantine-ledger figures), \texttt{qcos\_nonintegrable.py} (the
exact linear algebra behind \S\ref{sec:galois}), together with the Rust tool
\texttt{tools/qcos\_rank} for the rank exclusions of Remark~\ref{rem:zerocurve-arith} and
\texttt{beta2\_effective\_irrationality.py} for Proposition~\ref{prop:effective}.

\section{The $q$-difference module}\label{sec:prelim}

Throughout $Q=q^2$, $\sigma f(z)=f(Qz)$, and $a(z)=q+1-qz$. From the ${}_1\phi_1$ recurrence, $G$
satisfies
\begin{equation}\label{eq:Leq}
   q\,G(z)-a(z)\,G(Qz)+G(Q^2z)=0,\qquad L:=q-a\sigma+\sigma^2,
\end{equation}
a second-order $q$-difference equation with companion matrix $A(z)=\bigl(\begin{smallmatrix}0&1\\-q&
a\end{smallmatrix}\bigr)$, $\det A=q$, regular (Poincar\'e-type) at $z=0$ and irregular at $z=\infty$
(Newton slopes $\{0,1\}$). A second solution is $\widetilde G(z)=z^{1/2}H(z)$ with
$H(z)=\sum_{k\ge0}(-1)^kq^{k^2}(1-q)z^k/(q;q)_{2k+1}$; the pair $(G,H)$ realizes the Hahn--Exton
$q$-cosine/$q$-sine $J^{(3)}_{\mp1/2}$. The Casoratian $\Cas(G,\widetilde G)(z):=G(z)\widetilde
G(Qz)-G(Qz)\widetilde G(z)=c\,z^{1/2}$ is nonzero, and the exact $q$-Wronskian identity
\begin{equation}\label{eq:qwron}
   q\,G(z)\,H(Qz)-G(Qz)\,H(z)=q-1
\end{equation}
holds; it is reproduced to the full working precision, at least $120$ digits, at three
$(q,z)$ points and at two precisions by \texttt{qcos\_analytic.py}. We write $M$ for the associated rank-$2$ $\sigma$-module over $\C(z)$.

At $z=\infty$ the equation has a formal solution basis reflecting the two Newton slopes. With
$\theta_Q(x)=\sum_{n\in\Z}Q^{n(n-1)/2}x^n$, these are
\[
   f_1(z)=\theta_Q(-z)\,\hat g(z),\quad \hat g(z)=\sum_{k\ge0}b_kz^{-k},\ \ b_k=\frac{q^{\,k-k(k-1)/2}}{(q;q)_k};
   \qquad
   f_2(z)=\frac{h(z)}{\theta_Q(-z/q)},\quad h(z)=\sum_{k\ge0}c_kz^{-k},
\]
where $\hat g$ is $q$-Gevrey divergent (its $q$-Borel transform is $1/(q^2/\xi;q)_\infty$, so
$f_1$ is Borel--Laplace summable off one Stokes direction) and $h$ is convergent ($c_k=O(Q^{k^2/2})$).
The Newton polygon, the $q$-Gevrey classification and the $q$-Borel--Laplace summation used here are
those of Ramis, Sauloy and Zhang \cite{RSZ}.
The coefficients $b_k$, which drive the Galois argument below, satisfy
$b_k\asymp q^{-k^2/2}$; both $\hat g$ and $h$ are determined by \eqref{eq:Leq} and $b_0=c_0=1$. Full
details of the resummation are in \cite{Bonfioli2}; here only the asymptotic expansions are used.

\section{The difference Galois group}\label{sec:galois}

\begin{lemma}[irreducibility of the $\sigma$-module]
\label{lem:irreducible}
For $0<q<1$ the rank-$2$ $\sigma$-module of $G$ is irreducible over $\C(z)$: the Riccati equation
\[
   u(z)\,u(Qz)-a(z)\,u(z)+q=0,\qquad a=q+1-qz,
\]
has no solution $u\in\C(z)$.
\end{lemma}

\begin{proof}
Suppose $u=P/R$ with $P,R\in\C[z]$ coprime, of degrees $p,r$, leading coefficients $\alpha,\beta$.
Clearing denominators,
\[
   F(z)\;:=\;P(Qz)P(z)-a(z)P(z)R(Qz)+q\,R(z)R(Qz)\;\equiv\;0 .
\]
\emph{At $z=0$:} $F(0)=P(0)^2-(q+1)P(0)R(0)+qR(0)^2=0$; if $R(0)=0$ then $P(0)=0$, contradicting
coprimality, and symmetrically; so $P(0)R(0)\ne0$ and $u(0)\in\{1,q\}$.
\emph{Top degree:} the three terms have degrees $2p$, $p+r+1$, $2r$. If $p=r$, the top degree
$2p+1$ is attained by the middle term alone, with coefficient $q\alpha\beta Q^{r}\ne0$;
if $|p-r|\ge2$, the top degree is attained by the first or third term alone, with coefficient
$\alpha^2Q^{p}$ or $q\beta^2Q^{r}$. Both are impossible, so $|p-r|=1$.
\emph{Root pairing:} let $\zeta\ne0$ be a root of $R$ of multiplicity $m$, so $P(\zeta)\ne0$. At
$z=\zeta/Q$ the three terms of $F$ have orders $k$, $k+m+[\,a(\zeta/Q)=0\,]$, $m_0+m$, where
$k=\operatorname{ord}_{\zeta/Q}P$ and $m_0=\operatorname{ord}_{\zeta/Q}R$; since $F\equiv0$ the
minimum order cannot be attained by a single term, forcing $k\ge m_0+m\ge m$. Symmetrically, a root
$\xi\ne0$ of $P$ of multiplicity $k$ forces $\operatorname{ord}_{Q\xi}R\ge k$ (coprimality gives
$\operatorname{ord}_{Q\xi}P=0$, so the first term has order $k$ at $\xi$, the third
$\operatorname{ord}_{Q\xi}R$, and the middle at least $k+\operatorname{ord}_{Q\xi}R$; again no
single minimum). The two maps $\zeta\mapsto\zeta/Q$, $\xi\mapsto Q\xi$ are mutually inverse on the
nonzero root sets, and chains close at length two by coprimality. Summing multiplicities,
$p\le r\le p$, so $p=r$: contradicting $|p-r|=1$.
\end{proof}

\begin{theorem}[the difference Galois group contains $\mathrm{SL}_2$]
\label{thm:sl2}
For every $0<q<1$, the difference Galois group over $\C(z)$ of the module of $G$ contains
$\mathrm{SL}_2$.
\end{theorem}

\begin{proof}
For $N\ge1$ suppose $\operatorname{Sym}^N\!M$ has a rank-$1$ submodule, i.e.\ there is a nonzero
solution $w$ of the $N$-th symmetric-power equation with $r:=w(Q\,\cdot)/w\in\C(z)$. The products
$f_1^{N-k}f_2^{\,k}$ ($0\le k\le N$) form a solution basis (their Casoratian is a power
of $\operatorname{Cas}(f_1,f_2)\ne0$, by the standard symmetric-power formula), so
$w=\sum_kc_kf_1^{N-k}f_2^k$ with $c_k$ in the field of $\sigma$-constants; set
$k_0:=\min\{k:c_k\ne0\}$. In the analytic setting those $\sigma$-constants are the
$Q$-elliptic functions rather than $\C$, which does not affect what follows: only
$c_{k_0}\ne0$ is used.

\emph{Case $k_0<N$.} Along a ray in a regular class, $f_2/f_1$ decays super-polynomially, so the
rational function $r$ has the Poincar\'e asymptotic expansion of the quotient of the dominant
component:
\[
   r\;\sim\;(-1/z)^{N-k_0}(-z/q)^{k_0}\,
   \Bigl(\tfrac{\hat g(Qz)}{\hat g(z)}\Bigr)^{N-k_0}
   \Bigl(\tfrac{h(Qz)}{h(z)}\Bigr)^{k_0}.
\]
The $j$-th coefficient of the formal series on the right is
$(N-k_0)\,b_jQ^{-j}\bigl(1+O(q^{\,j-1})\bigr)$ with $b_j=q^{\,j-j(j-1)/2}/(q;q)_j$: the dominant term
comes from $b_jQ^{-j}$ in $\hat g(Qz)/\hat g(z)$, and every competing product is smaller by at least
$q^{\,j-1}$. That gap is not a numerical observation but an identity. For $1\le i\le j-1$,
\begin{equation}\label{eq:supermult}
   \frac{b_i\,b_{j-i}}{b_j}\;=\;q^{\,i(j-i)}\begin{bmatrix}j\\i\end{bmatrix}_q ,
\end{equation}
since the powers of $q$ contribute
$\tfrac12\bigl(j(j-1)-i(i-1)-(j-i)(j-i-1)\bigr)=i(j-i)$ and the Pochhammer factors assemble into the
$q$-binomial coefficient, a polynomial in $q$ with constant term $1$. As $i(j-i)\ge j-1$ on that
range, every split-off product carries at least one factor $q^{\,j-1}$ relative to $b_j$, uniformly
in $q$ and in $N,k_0$. (An earlier version of this proof rested on a numerical check at a single $j$
and a single $q$, and exhibited no~$c$. The exponent identity, the polynomiality of
$\bigl[\begin{smallmatrix}j\\i\end{smallmatrix}\bigr]_q$ and the bound $i(j-i)\ge j-1$ are now
verified exactly over $\Z$ for $j\le399$ by \texttt{qcos\_identities.py}, and machine-checked as
\texttt{supermult\_exponent} and \texttt{gap\_ge}; see \S\ref{sec:lean}.) These coefficients grow like $q^{-j^2/2}$, whereas a rational
function's expansion at $\infty$ has geometrically bounded coefficients. Since matching asymptotic
expansions are unique, no rational $r$ exists.

\emph{Case $k_0=N$.} Then $r=(-z/q)^N\rho^N$ with $\rho:=h(Q\,\cdot)/h$ meromorphic on $\C^*$;
$\rho^N\in\C(z)$ and single-valuedness force $\rho\in\C(z)$ (an $N$-th root of a rational function
is single-valued on $\C^*$ only if every divisor multiplicity, including at $0$ and $\infty$, is
divisible by $N$). Then $f_2(Qz)/f_2(z)=-(z/q)\rho(z)$ is rational, giving a rational solution of
the rank-$2$ Riccati equation: contradicting Lemma~\ref{lem:irreducible}.

Hence no $\operatorname{Sym}^N\!M$ has a rank-$1$ submodule. $N=1$ is irreducibility; $N=2$
excludes the imprimitive (dihedral) case; and a projectively finite irreducible subgroup
($A_4,S_4,A_5$) fixes a line in some $\operatorname{Sym}^N$ with $N\le12$ (the classical binary
polyhedral invariants), which is likewise excluded. By the classification of algebraic subgroups of
$\mathrm{GL}_2$, as used in difference Galois theory in van der Put--Singer
\cite[Ch.~1]{vdPS}, the group contains $\mathrm{SL}_2$.
\end{proof}

The rank-two criterion used is that of Dreyfus, Hardouin and Roques
\cite[Prop.~2.9]{DHR2018}; the corresponding dichotomy for general second-order linear difference
equations, including the elliptic case, is Arreche, Dreyfus and Roques \cite{ADR}, which is later
and does not record \cite{DHR2018}.

\begin{proposition}[no $\sigma\delta$-integrability]
\label{prop:nonintegrable}
For every $0<q<1$ there is no $B\in\mathrm{gl}_2(\C(z))$ and no $c\in\C(z)$ with
$\theta(A)=\sigma(B)A-AB+cA$, where $A=\bigl(\begin{smallmatrix}0&1\\-q&a\end{smallmatrix}\bigr)$ and
$\theta=z\partial_z$.
\end{proposition}

\begin{proof}
Writing $B=\bigl(\begin{smallmatrix}\alpha&\beta\\\gamma&\delta\end{smallmatrix}\bigr)$, the entries give
$\gamma=-q\beta_1$, $\delta=\alpha_1+a\beta_1-c$, and eliminating $\alpha$ (the central terms $c$
cancel identically; verified symbolically in \texttt{qcos\_nonintegrable.py}) leaves the
necessary condition
\[
   qa\,\beta_3\;=\;a_1(aa_1-q)\,\beta_2\;-\;a(aa_1-q)\,\beta_1\;+\;qa_1\,\beta\;+\;q\,z\bigl(a_1+q^2a\bigr),
\]
$\beta_k=\beta(Q^kz)$, $a=a(z)$, $a_1=a(Qz)$. A polynomial part of degree $m\ge0$ contributes top
coefficient $\propto\beta_m(1-q^{2+2m})\ne0$ at degree $m+3$, above the inhomogeneity: so $\beta$
has no polynomial part. If $\beta$ has a pole off the orbit $z_aQ^{\Z}$ and away from $z=0$, the maximal-modulus pole
appears in exactly one term (the slot $z=\zeta Q^{-3}$, coefficient $qa\ne0$ there), and the
minimal-modulus pole likewise (slot $z=\zeta'$, coefficient $qa_1$): both impossible. The origin
must be treated separately, since at $\zeta=0$ one has $\zeta Q^{-3}=0$ as well and all four
$\beta$-terms are singular at once, so the isolation argument does not apply. Suppose then
$\beta\sim Bz^{-n}$ at $0$ with $B\ne0$ and $n\ge1$. Collecting the coefficient of $z^{-n}$ in the
displayed relation gives
\[
   (q+1)\,B\,q\,(x-1)(x-q)(x-1/q),\qquad x:=q^{-2n},
\]
which is nonzero for every $n\ge1$ because $x=q^{-2n}>1/q>1>q$; so $\beta$ is regular at the origin.
With that case closed, so the poles
lie in the window $\{z_aQ^{-1},\dots,z_aQ^{3}\}$ ($a(z_aQ^k)=(q+1)(1-Q^k)$ vanishes only at $k=0$).
Sweeping the slots $z=z_aQ^{-4},\dots,z_aQ^{-1}$ kills the orders at $z_aQ^{-1},\dots,z_aQ^{2}$ one
by one (each pole appears in a single term with nonvanishing coefficient), and the order at
$z_aQ^{3}$ is killed at the slot $z=z_aQ$ (coefficient $a_1(aa_1-q)$ with
$(aa_1-q)(z_aQ)=(q+1)^2(1-q^2)(1-q^4)-q$) or, where that vanishes (one root
$q\approx0.78027$), at the slot $z=z_aQ^{2}$: coefficient $a(aa_1-q)$ with
$(aa_1-q)(z_aQ^2)=(q+1)^2(1-q^4)(1-q^6)-q$ (one root $q\approx0.86753$); the two polynomials are
coprime, so for every $q\in(0,1)$ at least one slot applies. Hence $\beta\equiv0$, contradicting
$z(a_1+q^2a)=z\bigl[(q+1)(1+q^2)-2q^3z\bigr]\not\equiv0$. (Cross-check, \texttt{qcos\_nonintegrable.py}: the exact linear
system for a $\beta$ with poles of order $\le3$ at the five window points has $15$ unknowns and
$\operatorname{rank}(A)=15$, $\operatorname{rank}[A|b]=16$, hence no solution; the same holds at
$q=\tfrac15$. The two slot polynomials are located there as well, with roots
$0.7802652115\ldots$ and $0.8675323943\ldots$ in $(0,1)$, and shown to be coprime over $\Q[q]$.)
\end{proof}

\begin{corollary}[differential transcendence of $G$]
\label{cor:hypertranscendence}
For every $0<q<1$ the function $G(q,\cdot)$ satisfies no algebraic differential equation over
$\C(z)$: the Hahn--Exton $q$-cosine is differentially transcendental.
\end{corollary}

\begin{proof}[Proof, via parametrized difference-Galois theory]
We apply the parametrized ($\sigma,\theta$)-difference Galois theory with derivation
$\theta=z\partial_z$. The determinant equation $\sigma(y)=\det(A)\,y=q\,y$ has the solution
$y=z^{1/2}$ (indeed $\sigma(z^{1/2})=(Qz)^{1/2}=q\,z^{1/2}$, as $Q=q^2$), which is algebraic over $\C(z)$; hence the parametrized difference Galois group of $\sigma(y)=\det(A)y$
lies in $\C^\times$. This is exactly the hypothesis of Dreyfus--Hardouin--Roques
\cite[Prop.~2.9]{DHR2018}, which for a rank-$2$ system whose difference Galois group contains
$\mathrm{SL}_2$ (here Theorem~\ref{thm:sl2}) gives the dichotomy: the parametrized group
$\mathrm{Gal}^\theta$ either is conjugate into $\mathrm{GL}_2(\C)$ (all solutions differentially
algebraic) or contains $\mathrm{SL}_2(\widetilde{\mathcal C})$ (all nonzero solutions
differentially transcendental), and the \emph{first} alternative holds if and only if there exists
$B\in\mathrm{gl}_2(\C(z))$ with $\theta(A)=\sigma(B)A-AB$ (their integrability condition~(7);
$\theta(\det A)=0$ since $\det A=q$ is constant in $z$, so no trace term appears). By
Proposition~\ref{prop:nonintegrable} no such $B$ exists: indeed none exists even allowing a
scalar twist $cA$. Therefore the second alternative holds:
$\mathrm{Gal}^\theta\supseteq\mathrm{SL}_2(\widetilde{\mathcal C})$, and every nonzero solution,
in particular $G(q,\cdot)$, is differentially transcendental over $\C(z)$.
\end{proof}


\section{The connection constant and the zero product}\label{sec:zeroproduct}

The behaviour of $G$ at $z=\infty$ is governed by a theta-comparison. Let
$\theta_Q(x)=\sum_{n\in\Z}Q^{n(n-1)/2}x^n$, so $\theta_Q(Qx)=x^{-1}\theta_Q(x)$ and, by the Jacobi
triple product \cite{Jacobi1829} (see \cite[(1.6.1)]{GR}),
$\theta_Q(-z)=(Q;Q)_\infty(z;Q)_\infty(Q/z;Q)_\infty$, an entire function of order $0$
with simple zeros exactly at $z\in Q^{\Z}$. The following theorem evaluates the leading asymptotic
constant of $G$ and, as a consequence, the regularized product of its zeros; the constant
$C_1:=\lim_{R\to\infty}G(-R)/\theta_Q(R)$ coincides with the Stokes-invariant connection coefficient of
$G$ against the theta solution.

\begin{theorem}[the connection constant, and the zero product of the Hahn--Exton $q$-cosine]
\label{thm:zeroproduct}
For $0<q<1$,
\[
   C_1=\frac1{(q;q)_\infty},
   \qquad\text{and consequently}\qquad
   \prod_{k\ge1} z_k(q)\,Q^{\,k-1}\;=\;(q;q^2)_\infty ,
\]
where $z_1<z_2<\cdots$ are the positive zeros of $G(q,\cdot)$. Likewise, for the second solution's
underlying entire function $H$ with zeros $w_1<w_2<\cdots$,
$\prod_{k\ge1}w_k(q)\,q\,Q^{\,k-1}=(q^3;q^2)_\infty$. \textup{(\texttt{qcos\_analytic.py}:
both products agree with the right-hand sides to the full working precision, $120$ and $240$
digits, at the three bases $q=0.3$, $q=0.42$ and $q=q^\ast$; the constant $C_1$ is approached
along $z=uQ^{-N}$ at the rate $O(q^N)$, with $|(q;q)_\infty G/\theta_Q(-z)-1|=3.7\cdot10^{-48}$
at $q=0.3$, $N=45$.)}
\end{theorem}

\noindent An earlier version of this note attributed the constant $C_1=1/(q;q)_\infty$ to
Morita, as the specialisation to the $q$-cosine ($\nu=-\tfrac12$) of his connection formula for
$J^{(3)}_\nu$ \cite[Thm.~1]{Morita2011}. That attribution is withdrawn, on three counts. The base
dictionary was wrong: Morita works with base $p$ and $q=p^2$, so his $p$ is the present $q$, not
$q^2$. His Theorem~1 carries the coefficients $1/(p^{\mp2\nu},p;p)_\infty$, and
$(p^{-1};p)_\infty=0$, so the formula is \emph{singular} at exactly $2\nu=\pm1$, which is the
$q$-cosine, the case wanted here. And he expands the solution at infinity in the basis at the
origin, whereas $C_1$ is the coefficient against the divergent theta solution, which does not
appear in his Theorem~1 at all. The constant is therefore proved here, not quoted, and the
regularised zero product is a corollary of it rather than of \cite{Morita2011}. Step~(3) of the proof below re-derives the zero localisation
$z_k=Q^{1-k}(1+O(q^k))$ of Abreu--Bustoz--Cardoso \cite[Thm.~2.3]{ABC2003}.

\begin{proof}
\emph{(1) An exact identity.} From $(q;q)_{2k}=(q;q)_\infty/(q^{2k+1};q)_\infty$ and Euler's
expansion $(x;q)_\infty=\sum_{j\ge0}(-1)^jq^{j(j-1)/2}x^j/(q;q)_j$ \cite[(1.3.16)]{GR} at
$x=q^{2k+1}$, together with
$k(k-1)+2kj=(k+j)(k+j-1)-j(j-1)$ and the substitution $m=k+j$, absolute convergence of the double
sum permits the rearrangement
\[
   (q;q)_\infty\,G(z)\;=\;\sum_{m\ge0}(-1)^mq^{m(m-1)}z^m\,P_m(1/z),
   \qquad
   P_m(w):=\sum_{j=0}^m c_jw^j,\quad c_j=\frac{q^{\,j-j(j-1)/2}}{(q;q)_j}
\]
(the $c_j$ coincide with the $b_j$ of \S\ref{sec:prelim}; the identity is checked in
$\Q[[q]]$, coefficient of $z^k$ by coefficient of $z^k$, to order $q^{90}$ or better for $k\le6$,
and numerically to the full working precision at two precisions:
\texttt{qcos\_identities.py}).

\emph{(2) The constant.} Fix $u<0$ and put $z=uQ^{-N}$. The moduli $q^{m(m-1)}|z|^m$ peak at
$m_0=N+O(1)$ and decay like $q^{(m-m_0)^2+O(1)}$ away from it. For $1\le j\le m\le 2N$,
$2Nj-j(j-1)/2\ge Nj$, so $|c_jz^{-j}|\le Cq^{Nj}|u|^{-j}$ and
$|P_m(1/z)-1|\le C'q^{N}$ uniformly; the terms with $m>2N$ contribute a relative
$O(q^{N^2})$. The negative-$m$ part of $\theta_Q(-z)$ is $O(|z|^{-1})$, negligible against the peak.
Hence $(q;q)_\infty G(z)/\theta_Q(-z)\to1$ along $z=uQ^{-N}$, for each fixed $u<0$.

\emph{(3) Geometric zero dressing.} On the circle $|z-Q^{1-k}|=\varepsilon q^kQ^{1-k}$ the identity
of step (1) and the bounds of step (2) give
$|(q;q)_\infty G-\theta_Q(-z)|\le Cq^k|\theta_Q(-z)|_{\text{peak}}$, while $\theta_Q(-z)$ has a
simple zero at $Q^{1-k}$ with derivative bounded below there; Rouch\'e gives
$z_k=Q^{1-k}(1+O(q^k))$, i.e.\ $\log(z_kQ^{k-1})=O(q^k)$, absolutely summable.

\emph{(4) The product.} $G(z)=\prod_k(1-z/z_k)$ (order $0$, genus $0$, $G(0)=1$) and
$\theta_Q(-z)=(Q;Q)_\infty(z;Q)_\infty(Q/z;Q)_\infty$ (Jacobi). Along $z=-R$, $R\to\infty$,
\[
   \frac{G(-R)}{\theta_Q(R)}
   =\frac1{(Q;Q)_\infty(-Q/R;Q)_\infty}\prod_{k\ge1}\frac{1+R/z_k}{1+RQ^{k-1}} ,
\]
and for $x,y>0$ the function $R\mapsto\log\frac{1+Rx}{1+Ry}$ is monotone from $0$ to $\log(x/y)$, so
each factor is dominated by $|\log(z_kQ^{k-1})|$, summable by step (3). Dominated convergence gives
$G(-R)/\theta_Q(R)\to1/\bigl[(Q;Q)_\infty\prod_kz_kQ^{k-1}\bigr]$ as $R\to\infty$; evaluating along
the subsequence $R=|u|Q^{-N}$ by step (2), the limit equals $1/(q;q)_\infty$, whence
$C_1=1/(q;q)_\infty$ and $\prod_kz_kQ^{k-1}=(q;q)_\infty/(Q;Q)_\infty=(q;q^2)_\infty$ by Euler's
$(q;q)_\infty=(q;q^2)_\infty(q^2;q^2)_\infty$. The $H$-statement is identical with comparison
function $\theta_Q(-qz)$ and ratio $(1-q)/(q;q)_\infty$.
\end{proof}

\begin{remark}[the modular split at $z^\ast$]
\label{rem:zeroproduct-dioph}
At $q=q^\ast$ the product splits off the candidate rational zero:
\[
   \prod_{k\ge2}z_k(q^\ast)\,Q^{\,k-1}=\frac{(q;q^2)_\infty}{z^\ast}
   =0.9882998846287462938546394\ldots
\]
Were $z^\ast$ (equivalently $\beta_2$) algebraic, the modular value $(q;q^2)_\infty$ (transcendental at algebraic $q$ by Nesterenko's theorem \cite{Nesterenko} applied to the Euler
products) would equal an algebraic number times the tail product $\prod_{k\ge2}z_kQ^{k-1}$. This is
consistent only if the tail product is itself transcendental; the arithmetic of the zero condition is
thereby confined to that tail, which is the content of the multiplicity problem of \cite{Bonfioli2}.
\end{remark}


\section{The arithmetic of the zeros}\label{sec:arith}

\subsection{Conjecture Z and its confluent evidence}
Conjecture~\ref{conj:Z} is the $q$-analogue of the classical theorem (Siegel; see also Lindemann for
the endpoint) that the positive zeros $j_{\nu,k}$ of the Bessel function $J_\nu$ are transcendental for
rational $\nu$. Two structural facts support it. First, Theorem~\ref{thm:main-transc}: $G$ is
differentially transcendental, so its zeros are not the zeros of any algebraic-differential shadow
whose arithmetic could be controlled classically. Second, in the confluent limit $q\to1$ the equation
\eqref{eq:Leq} degenerates to Bessel's equation, with the zeros rescaling as
\[
   \frac{z_k(q)}{(1-q)^2}\;\longrightarrow\;j_{-1/2,k}^{\,2}
   \qquad(q\to1^-),
\]
so Conjecture~\ref{conj:Z} deforms the classical statement. The scaling factor matters and an
earlier version omitted it, writing $z_k(q)\to(qj_{-1/2,k})^2$; that is false, the zeros collapsing
to $0$ rather than to a finite limit. At $k=1$ the measured ratios $z_1/(1-q)^2$ are
$2.33902$, $2.40449$, $2.45502$ at $q=0.9,0.95,0.99$ against $j_{-1/2,1}^2=(\pi/2)^2=2.46740$,
whereas $z_1(0.99)=2.45502\cdot10^{-4}$ against the discarded $(q\pi/2)^2=2.4183$. Here $z_1$ is
located directly as the least positive zero of $G(q,\cdot)$ and not through the series of
Theorem~\ref{thm:integrality}, which at these bases is outside its disc of convergence
(\texttt{qcos\_analytic.py}).

That confluence also explains, in one line, why the classical case is provable and the $q$-case is
not. Write $D_N$ for the denominator of the $N$-th partial sum. The classical engine (Lambert,
Lindemann, and every factorial-denominator proof) runs on a \emph{divergent step ratio}:
$D_{N+1}/D_N=N+1\to\infty$, so one further term is gained at asymptotically no arithmetic cost and
$\text{tail}\cdot D_N\asymp2/(N+1)\to0$ (measured: $0.3097,\,0.1629,\,0.0991$ at $N=3,6,10$ for
$e=\sum1/k!$). The $q$-deformation replaces the factorial by $(q;q)_k$, whose step ratio is
\emph{constant} in the relevant scale: each further term gains $\asymp q^{2N}$ but costs
$\asymp t^{2N}$, and at $s=1$ the two cancel identically, so $\text{tail}\cdot D_N$ diverges
(measured at $q=\tfrac13$, $z_0=2q(1-q)$, with $D_N=t^{\deg_qP_N}=3^{2N^2+N}$:
$7.15,\,7.47\cdot10^5,\,5.14\cdot10^{14}$ at $N=2,4,6$; including the argument denominator $9^N$
only makes it worse). The bridge is explicit:
$(q;q)_k/\bigl((1-q)^kk!\bigr)\to1$ as $q\to1$ (measured at $k=6$: $0.4719,\,0.9277,\,0.9925$ at
$q=0.9,\,0.99,\,0.999$). All three tables are in \texttt{qcos\_ledger.py}. The factorial (and with it the divergent step ratio that powers
Lindemann) exists only in the confluent limit; at any $q<1$ the ratio is constant, and a constant
ratio can never win. This is the same exactness of Remark~\ref{rem:thetabarrier} seen dynamically,
and it identifies precisely what a proof of Conjecture~\ref{conj:Z} must supply: an arithmetic
mechanism with no classical shadow.

\subsection{The self-referential specialization: $\beta_2$}
The reflection-group growth constant $\beta_2$ is the reciprocal square root of the least positive root
$q^\ast$ of $S(q)=G(q,2q(1-q))$, which is exactly the condition that the catalytic point
$z^\ast=2q^\ast(1-q^\ast)$ be a zero of $G(q^\ast,\cdot)$. Thus $\beta_2$ is governed by a
\emph{self-referential} specialization of Conjecture~\ref{conj:Z}: the zero and the base move together
along the diagonal $z=2q(1-q)$.

At this point the $q$-Wronskian \eqref{eq:qwron} yields a clean exact relation between two Hahn--Exton
values: since $G(q^\ast,z^\ast)=0$,
\begin{equation}\label{eq:wron-zero}
   G(q^\ast,\,q^{\ast2}z^\ast)\cdot H(q^\ast,\,z^\ast)=1-q^\ast
\end{equation}
(reproduced to the full working precision at two precisions, at least $120$ digits,
by \texttt{qcos\_analytic.py}): a product of two $q$-Bessel values at the algebraic point $z^\ast$ equal
to the algebraic number $1-q^\ast$. Together with $G(q^\ast,z^\ast)=0$ and the modular zero product of
Theorem~\ref{thm:zeroproduct-intro}, this places three exact constraints on the values at $z^\ast$; a
transcendence proof requires one more independent constraint than the reductio supplies, which is the
precise sense in which Conjecture~\ref{conj:Z} is a multiplicity (rather than a single-relation)
problem.

\subsection{The $q$-Siegel--Shidlovskii margin}
The natural route to Conjecture~\ref{conj:Z} is the $q$-analogue of the Siegel--Shidlovskii method
developed by Amou--Matala-aho--V\"a\"an\"anen \cite{AMV}. Two things must be said precisely, since
the terminology is ours and not theirs. Their theorem applies to systems of the shape
$z^sy(qz)=a(z)Cy(z)+b(z)$ with $C$ a \emph{constant} matrix, not to Poincar\'e-type systems in
general; and their hypothesis is arithmetic on $q$, namely (5.4) of \cite{AMV} in terms of
$\lambda=\log H(q)/\log\|q\|_v$, not a margin in our sense. What their theorem yields is
\emph{linear independence over the fixed field} $\mathbb K$ containing $q$, never transcendence over
$\overline{\Q}$; no $q$-analogue of Siegel--Shidlovskii giving transcendence is known. The
quantity we call the margin below is a convenient bookkeeping ratio for the same obstruction: the
ratio of the logarithmic height of the cleared denominators to that of the analytic decay. A reader
working on the value side should compare Amou--Katsurada--V\"a\"an\"anen \cite{AKV2002}, who for a
$q$-analogue of $J_0$ determine exactly which $\alpha\in\mathbb K^\times$ have value in $\mathbb K$;
we make no value-level claim of that kind here, and the question addressed below concerns the
\emph{zeros}, on which we have found nothing in the literature. For the Hahn--Exton class this ratio is
\[
   \gamma=\frac{\log(t/s)}{2\log t}\le\frac12\qquad(q=s/t\text{ in lowest terms}),
\]
the factor $2$ being the double Pochhammer $(q;q)_{2k}$; this is below the threshold, so the method does
not apply as stated. Extending it past the margin for this irregular-at-$\infty$ class is the
outstanding problem; the connection data needed for the irregular direction is available in closed form
\cite{Bonfioli2}.

\subsection{Integrality of the zero lattice}
The zeros, viewed as functions of the base, are integer power series after the Newton-polygon
rescale: all of them, not only the first.

\begin{theorem}[integrality of the zero lattice]\label{thm:integrality}
For every $k\ge1$ the rescaled $k$-th zero $u_k(q):=q^{2(k-1)}z_k(q)$ lies in $1+q\,\Z[[q]]$.
In particular
\[
   z_1(q)=1-q-q^3+q^4-q^5+2q^6-2q^7+4q^8-6q^9+8q^{10}-14q^{11}+21q^{12}-\cdots\ \in\ \Z[[q]] .
\]
\end{theorem}

\noindent The \emph{shape} of such an expansion is known: Ismail and Zhang \cite{IsmailZhang2007}
show that $q^{2n-1}z_n$ is analytic in $q^n$, with coefficients polynomial in three generators over
rational functions of $q$, and the relevant entire function lies in the admissible class of
Hayman \cite{Hayman1956}.
What is proved here, and what we have not found in that literature, is that the coefficients are
\emph{integers}.

\begin{proof}
From $g_{k'}=(-1)^{k'}q^{k'(k'-1)}/(q;q)_{2k'}$ and the exponent identity
$k'(k'-1)-2(k-1)k'+k(k-1)=(k'-k)(k'-k+1)$,
\[
   F_k(q,u):=q^{k(k-1)}\,G\bigl(q,\,u\,q^{-2(k-1)}\bigr)
   =\sum_{k'\ge0}(-1)^{k'}q^{(k'-k)(k'-k+1)}\frac{u^{k'}}{(q;q)_{2k'}} .
\]
The exponent $(k'-k)(k'-k+1)$ is a product of consecutive integers, hence nonnegative, and vanishes
exactly at $k'\in\{k-1,k\}$: each zero scale is carried by its own edge of the strictly convex
Newton polygon of $G$. Since $1/(q;q)_{2k'}\in\Z[[q]]$ (partitions into parts $\le2k'$),
$F_k\in\Z[[q]][[u]]$ with $F_k(0,u)=(-1)^{k-1}u^{k-1}(1-u)$; at $u=1$ the value vanishes and
$\partial_uF_k(0,1)=(-1)^k$ is a unit, so Hensel's lemma in the complete local ring $\Z[[q]]$
yields a unique root $u_k\in1+q\,\Z[[q]]$. The identification of this formal root with the
analytic $k$-th zero is not formal: it holds where the series converges, and $u_1$ has radius
$|q_c|=0.5546\ldots$ by Remark~\ref{rem:zerocurve-arith}(e). For $0<q<|q_c|$ the scale
localization $z_k=q^{2(1-k)}(1+O(q^k))$ of \S\ref{sec:zeroproduct} pins the root to the $k$-th
zero; outside that disc the statement of the theorem is about the power series and not about the
zero. (The lattice is regenerated exactly in $\Z[[q]]$ to order $q^{300}$ for $k\le8$ by
\texttt{qcos\_lattice.py}: $u_2=1-q^6-2q^8+q^9-\cdots$, $u_3=1-q^{15}-\cdots$,
$u_4=1-q^{28}-\cdots$. The onset law visible here is proved next. The Newton recursion is carried
out over $\Z$ and its residual verified to order $q^{80}$ for $k\le6$ in Lean; the Hensel step
itself is not formalised, see \S\ref{sec:lean}.)
\end{proof}

\subsection{Lattice values and the hexagonal law}
The deviation $u_k-1$ begins exactly at the $k$-th hexagonal number $k(2k-1)$. This reduces to, and
is proved by, an exact evaluation of $G$ at its own lattice points.

\begin{lemma}[parameter--argument swap; Koornwinder--Swarttouw]\label{lem:swap}
For $|Q|<1$,
$(b;Q)_\infty\,{}_1\phi_1(0;b;Q,z)$ is symmetric in $(b,z)$; equivalently
\begin{equation}\label{eq:dressing}
   (q;Q)_\infty\,G(q,z)=\sum_{n\ge0}\frac{(-1)^nq^{n^2}}{(Q;Q)_n}\,(zQ^n;Q)_\infty .
\end{equation}
\end{lemma}

\noindent This is Koornwinder--Swarttouw \cite[Prop.~2.1, eq.~(2.3)]{KS1992}, with the same double-sum
proof; they note in their Remark~2.2 that it is a limit case of Heine's transformation
\cite[\S1.4]{GR}. It is restated
here for use below, not claimed.

\begin{proof}
In $(b;Q)_\infty\,{}_1\phi_1(0;b;Q,z)=\sum_n(-1)^nQ^{\binom n2}z^n\,(bQ^n;Q)_\infty/(Q;Q)_n$
expand $(bQ^n;Q)_\infty=\sum_i(-1)^iQ^{\binom i2}(bQ^n)^i/(Q;Q)_i$ (Euler, \cite[(1.3.16)]{GR});
since
$\binom n2+\binom i2+ni=\binom{n+i}2$, the double sum is
$\sum_{n,i}(-1)^{n+i}Q^{\binom{n+i}2}z^nb^i/\bigl((Q;Q)_n(Q;Q)_i\bigr)$, manifestly symmetric.
For \eqref{eq:dressing} take $b=q$ and use $(z;Q)_\infty/(z;Q)_n=(zQ^n;Q)_\infty$ on the swapped
side, with $Q^{\binom n2}q^n=q^{n^2}$. (\texttt{qcos\_identities.py}: the symmetry holds to
the full working precision at a generic $(Q,b,z)$ at two precisions, and \eqref{eq:dressing}
itself is verified exactly in $\Q[[q]]$ to order $q^{250}$ at three rational $z$.)
\end{proof}

\begin{proposition}[square vanishing at the lattice]\label{prop:latticevalues}
For every $m\ge0$,
\[
   G\bigl(q,Q^{-m}\bigr)
   =(-1)^{m+1}\,\frac{(Q;Q)_\infty}{(q;Q)_\infty}
    \sum_{r\ge0}\frac{(-1)^r\,q^{(m+1+r)^2}}{(Q;Q)_r\,(Q;Q)_{m+1+r}} ,
\]
so $\operatorname{ord}_qG(q,Q^{-m})=(m+1)^2$ exactly, with leading coefficient $(-1)^{m+1}$.
\end{proposition}

\noindent This follows in two lines from Koornwinder--Swarttouw \cite[Prop.~2.1 and Rem.~2.3,
eq.~(2.6)]{KS1992}: putting $n=m+1$, $z=q$ in their (2.6) gives the leading factor $q^{(m+1)^2}$ and
the $r$-th term exponent $(m+1+r)^2$, which is the display above. It is stated separately only
because \S5 uses it repeatedly.

\begin{proof}
Evaluate \eqref{eq:dressing} at $z=Q^{-m}$. The factor $(Q^{n-m};Q)_\infty$ vanishes for $n\le m$
and equals $(Q;Q)_\infty/(Q;Q)_{n-m-1}$ for $n\ge m+1$; reindexing $n=m+1+r$ gives the display.
The $q$-order of the $r$-th term is $(m+1+r)^2$, strictly increasing in $r$, and all other factors
are units in $\Z[[q]]$, so no cancellation can occur and the $r=0$ term decides.
(\texttt{qcos\_identities.py} verifies the identity in $\Q[[q]]$ to order $q^{230}$ or better
for $m\le4$, confirms $\ord_qG(q,Q^{-m})=(m+1)^2$ there, and reproduces it numerically at $q=0.3$,
$m=3$ to $46$ and $107$ digits at working precisions $60$ and $120$.)
\end{proof}

\begin{corollary}[hexagonal law]\label{cor:hexagonal}
For every $k\ge1$,
\[
   u_k(q)=1-q^{\,k(2k-1)}\bigl(1+O(q)\bigr):
\]
the deviation of the $k$-th rescaled zero from $1$ begins exactly at the $k$-th hexagonal number,
with leading coefficient $-1$.
\end{corollary}

\begin{proof}
$F_k(q,1)=q^{k(k-1)}G(q,Q^{1-k})=(-1)^kq^{k(2k-1)}(1+O(q))$ by
Proposition~\ref{prop:latticevalues} at $m=k-1$ (note $k(k-1)+k^2=k(2k-1)$), while
$\partial_uF_k(q,1)=(-1)^k(1+O(q))$. The first Newton correction from $u=1$ is therefore
$-q^{k(2k-1)}(1+O(q))$, and all subsequent corrections are $O(q^{2k(2k-1)})$. This proves the
onsets $1,6,15,28$ observed above and predicts $u_5=1-q^{45}-\cdots$; the onsets
$k(2k-1)$ with leading coefficient $-1$ are confirmed for $k\le8$ in
\texttt{qcos\_lattice.py}.
\end{proof}

\begin{remark}\label{rem:latticevalues}
Three readings of Proposition~\ref{prop:latticevalues}. (a) The $r$-sum is again a Hahn--Exton
series, now of \emph{integer} order $m+1$ in base $Q$: the $q$-cosine ($\nu=-\tfrac12$) sampled
on its own zero lattice reproduces the integer-order family: an order reflection
$-\tfrac12\mapsto m+1$. (b) Along any backward orbit $z=Q^{-m}w$, \eqref{eq:dressing} gives
$G(Q^{-m}w)\sim(Q^{-m}w;Q)_\infty/(q;Q)_\infty$ as $m\to\infty$, with constant
$(w;Q)_\infty(Q/w;Q)_\infty/(q;Q)_\infty=\theta_Q(-w)/(q;q)_\infty$ by the triple product: recovering the connection constant of the backward normal form independently. (c) At any zero
$z_k(q)$ of $G$, \eqref{eq:dressing} transfers the vanishing into the \emph{parameter} slot:
${}_1\phi_1\bigl(0;z_k(q);Q,q\bigr)=0$. At the crossing $z_1(q^\ast)=2q^\ast(1-q^\ast)$ of the
reductio this is a ${}_1\phi_1$-vanishing with all-rational parameters; its truncation ledger has
the same heights as the original (the swap preserves the exactness of
Remark~\ref{rem:thetabarrier}).
\end{remark}

\begin{proposition}[sine side: integrality, lattice values, onsets]\label{prop:sinelattice}
All of the above has a sine-side counterpart. (a) The zeros $w_k(q)$ of $H$ satisfy
$v_k:=q^{2k-1}w_k(q)\in1+q\,\Z[[q]]$: with $(q^2;q)_{2k'}=(q;q)_{2k'+1}/(1-q)$,
\[
   q^{k(k-1)}\,H\bigl(q,\,v\,q^{1-2k}\bigr)
   =\sum_{k'\ge0}(-1)^{k'}q^{(k'-k)(k'-k+1)}\frac{v^{k'}}{(q^2;q)_{2k'}},
\]
the same strictly convex exponent as in Theorem~\ref{thm:integrality}, so Hensel applies on every
edge. (b) Writing $H(q,z)={}_1\phi_1(0;q^3;Q,qz)$, the dressing evaluates $H$ on its own lattice:
for $m\ge0$,
\[
   H\bigl(q,q^{1-2m}\bigr)=(-1)^m\,\frac{(Q;Q)_\infty}{q\,(q^3;Q)_\infty}
   \sum_{r\ge0}\frac{(-1)^r\,q^{(m+r+1)^2}}{(Q;Q)_r\,(Q;Q)_{m+r}},
\]
so $\operatorname{ord}_qH(q,q^{1-2m})=(m+1)^2-1$. (c) Consequently
$v_k=1-q^{\,k(2k+1)}(1+O(q))$ for every $k\ge1$.
\end{proposition}

\begin{proof}
(a) is the displayed identity together with the argument of
Theorem~\ref{thm:integrality}: the exponent $(k'-k)(k'-k+1)$ is strictly convex in $k'$ with a
unique minimiser, so the Newton polygon has a single edge of slope $0$ and Hensel applies.

(b) Apply Lemma~\ref{lem:swap} with $b=q^3$ and $Q=q^2$ to $H(q,z)={}_1\phi_1(0;q^3;Q,qz)$:
\[
   (q^3;Q)_\infty\,H(q,z)\;=\;\sum_{n\ge0}(-1)^n q^{\,n^2+2n}\,
      \frac{(qzQ^n;Q)_\infty}{(Q;Q)_n}.
\]
At $z=q^{1-2m}$ one has $qz=q^{2-2m}=Q^{1-m}$, so $qzQ^n=Q^{n+1-m}$ and the factor
$(Q^{n+1-m};Q)_\infty$ vanishes for $n\le m-1$, the product then containing the factor $1-Q^0$.
For $n\ge m$ it equals $(Q;Q)_\infty/(Q;Q)_{n-m}$. Writing $n=m+r$ and using
$(m+r)^2+2(m+r)=(m+r+1)^2-1$ turns the sum into the displayed one, the $q^{-1}$ prefactor being
the $-1$ in that exponent. The order statement follows since the $r=0$ term is the unique one of
minimal $q$-order.

(c) Combine (a) and (b): $v_k$ is the Hensel root attached to the edge, and the first correction
to $v_k=1$ enters at the $q$-order of the rescaled function at $v=1$, which by (b) at $m=k$ is
$k(k-1)+\operatorname{ord}_qH(q,q^{1-2k})=k(k-1)+\bigl((k+1)^2-1\bigr)=k(2k+1)$, exactly as on the
cosine side with $T_{2k-1}$ replaced by $T_{2k}$ (identity \texttt{sine\_onset} of
\S\ref{sec:lean}).
(\texttt{qcos\_lattice.py} verifies integrality of $v_k$ and the onsets $k(2k+1)$, with
leading coefficient $-1$, exactly in $\Z[[q]]$ to order $q^{300}$ for $k\le6$.)
\end{proof}

\begin{remark}[the triangular law]\label{rem:triangular}
$k(2k-1)=T_{2k-1}$ and $k(2k+1)=T_{2k}$ are consecutive triangular numbers, so the deviation
onsets of the two interleaved zero lattices are
\[
   \operatorname{onset}(u_k-1)=T_{2k-1},\qquad
   \operatorname{onset}(v_k-1)=T_{2k},
\]
with leading coefficient $-1$ in both families: between them, the cosine and sine lattices switch
on at every triangular number exactly once, alternating between the families. (Verified exactly in
$\Z[[q]]$ through $u_8$, i.e.\ $T_{15}=120$, and $v_6$, i.e.\ $T_{12}=78$;
\texttt{qcos\_lattice.py}.)
\end{remark}

\begin{proposition}[stable deviation law]\label{prop:stablelaw}
The two deviation series stabilize, $q$-adically as $k\to\infty$, to one and the same
eta-quotient:
\[
   \bigl(1-u_k(q)\bigr)q^{-T_{2k-1}}\ \longrightarrow\ \frac1{(q^2;q^2)_\infty^{2}}\
   \longleftarrow\ \bigl(1-v_k(q)\bigr)q^{-T_{2k}} ,
\]
whose coefficients $1,2,5,10,20,36,65,\dots$ count two-colored partitions. The agreement depth
measured exactly in $\Z[[q]]$ is $2k-2$ on the cosine side for $2\le k\le8$ and $2k-1$ on the
sine side for $2\le k\le6$ (\texttt{qcos\_lattice.py}).
\end{proposition}

\begin{proof}
The convergence asserted is $q$-adic (coefficient-by-coefficient in $\Z[[q]]$), so it suffices to
show each of three exact power-series limits below holds with agreement depth $\to\infty$ in $k$;
no analytic interchange is involved.

\emph{Newton correction.} Since $u_k$ is the Hensel root of $F_k(q,\cdot)$ at $u=1$ with
$F_k(q,1)=O(q^{T_{2k-1}})$ (Corollary~\ref{cor:hexagonal}) and $\partial_uF_k(q,1)$ a unit,
one Newton step gives
$1-u_k=F_k(q,1)/\partial_uF_k(q,1)\cdot(1+E_k)$ with $E_k=O(q^{T_{2k-1}})$ (the next correction is
quadratic in $F_k(q,1)$). As $T_{2k-1}=k(2k-1)\to\infty$, $E_k\to0$ $q$-adically, so it suffices to
evaluate the ratio $F_k(q,1)/\partial_uF_k(q,1)$ scaled by $q^{-T_{2k-1}}$.

\emph{Numerator.} By Proposition~\ref{prop:latticevalues} at $m=k-1$,
\[
   (-1)^kF_k(q,1)\,q^{-T_{2k-1}}
   =\frac{(Q;Q)_\infty}{(q;Q)_\infty(Q;Q)_k}+O(q^{2k+1})
   \ \xrightarrow{\ q\text{-adic}\ }\ \frac1{(q;Q)_\infty},
\]
the $r$-th term of the lattice sum contributing, after the prefactor $q^{k(k-1)}$ and the
rescaling by $q^{-T_{2k-1}}$, at order $q^{(k+r)^2-k^2}$, so that $r\ge1$ is $O(q^{2k+1})$, and
$(Q;Q)_k=(Q;Q)_\infty\bigl(1+O(q^{2k+2})\bigr)$. Agreement depth $2k+1$.

\emph{Denominator.} In $\partial_uF_k(q,1)=\sum_{k'\ge0}(-1)^{k'}k'\,q^{(k'-k)(k'-k+1)}/(q;q)_{2k'}$
the indices $k'=k+j$ and $k'=k-1-j$ share the exponent $(k'-k)(k'-k+1)=j(j+1)$ with opposite parity
and weights $k+j$, $k-1-j$; their sum has weight $(k+j)-(k-1-j)=2j+1$, whence
\[
   (-1)^k\partial_uF_k(q,1)
   =\frac{\sum_{j\ge0}(-1)^j(2j+1)q^{j(j+1)}}{(q;q)_\infty}+O(q^{2k-1})
   \ \xrightarrow{\ q\text{-adic}\ }\ \frac{(q^2;q^2)_\infty^{3}}{(q;q)_\infty}
\]
by Jacobi's identity $\sum_{j\ge0}(-1)^j(2j+1)q^{j(j+1)}=(q^2;q^2)_\infty^3$
\cite{Jacobi1829}. Agreement depth $2k-1$.

\emph{Quotient.} Dividing and using the exact splitting $(q;q)_\infty=(q;Q)_\infty(Q;Q)_\infty$,
\[
   \frac{1/(q;Q)_\infty}{(q^2;q^2)_\infty^{3}/(q;q)_\infty}
   =\frac{(q;q)_\infty}{(q;Q)_\infty(Q;Q)_\infty^{3}}
   =\frac{(Q;Q)_\infty}{(Q;Q)_\infty^{3}}
   =\frac1{(Q;Q)_\infty^{2}}=\frac1{(q^2;q^2)_\infty^{2}}.
\]
All three limits are exact with depth growing in $k$ and $E_k\to0$, so the product limit is exact;
this proves the cosine case. The sine case is identical with the numerator limit
$1/(q^3;Q)_\infty$ (Proposition~\ref{prop:sinelattice}) and $(q^2;q)_\infty=(q;q)_\infty/(1-q)$ in
the denominator, the two $(1-q)$ factors cancelling to give the same $1/(q^2;q^2)_\infty^{2}$.
(The measured cosine agreement depth is exactly $2k-2$ through $k=8$ and the sine exactly
$2k-1$ through $k=6$, consistent with the depth-$(2k-1)$ bound above.)
\end{proof}

\begin{remark}[arithmetic and analytic features of the zero curve]\label{rem:zerocurve-arith}
(a) The series satisfies no algebraic relation $P(q,z_1)=0$ of bidegree $\le(6,5)$; in particular
it is not a rational function of numerator and denominator degrees within that bidegree. Nothing
here excludes an algebraic relation of larger bidegree, and the statement is exactly the
bounded-profile one. \textup{(Rank certificate: the $42$ columns have full rank over the first
$106$ coefficients, at each of two primes; \texttt{tools/qcos\_rank}.)} (b) Its radius of
convergence is
$|q_c|=0.5545\ldots$, set by the real fold point of item (e); in particular
$q^\ast=0.4494\ldots$ lies inside the disc of convergence. (c) The
integrality opens a \emph{linear-height} approximation ledger: at $q=s/t$ the $N$-th partial sum is
rational with denominator $t^N$ and error $\asymp(q/R)^N$, $R=0.5545\ldots$, so the Liouville
criterion for the crossing $z_1(q)=2q(1-q)$ reads $s<R$: failing, but now only by the factor
$s/R\approx1.80$ at $s=1$, the narrowest margin among all ledgers recorded for this problem; the
obstruction has become the sub-unit radius, i.e.\ the complex collision of zeros. (d) The integer
coefficient sequence $1,-1,0,-1,1,-1,2,-2,4,-6,8,-14,21,-34,56,-88,148,-242,398,-669,1109,\dots$
appears to be new.

(e) The radius-limiting singularity is located and characterized: it is the simple fold
$q_c=-0.5545786101465797\ldots$, $z_c=2.5239525204417103\ldots$ (both regenerated to $230$
digits at two working precisions by \texttt{qcos\_analytic.py}), where the analytically continued branch collides with a second zero branch on the
negative $q$-axis ($G=G_z=0$ with $G_{zz}=0.2143\ldots\ne0$, $G_q=-1.5012\ldots\ne0$, so
$z_1(q)=z_c\pm\tilde C\sqrt{q-q_c}+\cdots$ with $\tilde C^2=2G_q/(-G_{zz})$ at the fold). The
fold law now has a closed-form constant:
\[
   |c_n|\,n^{3/2}|q_c|^n\ \longrightarrow\ \frac{|\tilde C|\sqrt{|q_c|}}{2\sqrt\pi}
   =0.7863083135\ldots,
\]
The constant is computed from the fold data at two working precisions, which agree to $230$
digits (\texttt{qcos\_analytic.py}). The sequence $|c_n|n^{3/2}|q_c|^n$ itself is strongly preasymptotic:
at $n=299$ it has reached only $0.78295$, agreeing with the limit to three digits, so the
coefficient asymptotics \emph{corroborate} $R=|q_c|$ rather than establish it. What does establish
$R=|q_c|$ is the fold itself, an actual singularity of the analytic continuation at that modulus.
Neither
$q_c$ nor $z_c$ satisfies an integer polynomial relation of degree $\le8$ with coefficients up to
$10^{14}$, nor does the pair satisfy a bilinear one (PSLQ at $230$ digits, tolerance $10^{-200}$).
Larger pair profiles are not reported: $230$ digits do not support more than about $16$ unknowns
against a $10^{14}$ coefficient bound, and beyond that PSLQ returns vectors that re-evaluation at
higher precision exposes as spurious.

(f) No change of variable rescues the ledger of (c): if $q=\psi(u)$ with $\psi\in u\,\Z[[u]]$ (as
required to keep $z_1\circ\psi\in\Z[[u]]$) and $u^\ast=\varphi(q^\ast)=s'/t'\in\Q$, the criterion
again reads $|u^\ast|\,t'=s'<R_u$; but a non-polynomial integer power series has radius at most $1$,
so $R_u\le1$ while $s'\ge1$. Every integral uniformization is capped at the boundary case
$s'=1=R_u$: the sub-unit radius can be moved but not beaten. This is the geometric face of the same
boundary on which the theta and Rogers--Ramanujan series sit in the truncation ledger
(Remark~\ref{rem:thetabarrier}).

(g) The cap has an elementary final form, the \emph{discreteness cap}: for any non-polynomial
$X\in\Z[[q]]$, any $q_0=s/t$ inside its disc of convergence (away from the boundary), and any $N$,
the truncation tail satisfies $|X(q_0)-X_N(q_0)|=|c_M|(s/t)^M(1+o(1))\ge s^M\,t^{-M}(1-o(1))$, where
$M>N$ is the first index with $c_M\ne0$: at least the Liouville floor $t^{-M}$ of the truncation's
denominator whenever $s\ge1$, because a nonzero integer coefficient has $|c_M|\ge1$. (Measured exactly on
$z_1$ at $q=\tfrac17$: tail/floor $=1.533,\,11.56,\,70.86,\,541.3$ at $N=6,10,14,18$;
\texttt{qcos\_ledger.py}.) Integrality thus
caps its own ledger: the property that grants linear heights also forbids the tail from undercutting
the floor. (h) The arithmetic class of $z_1$ is bounded from above as well. No linear differential
equation $\sum_{i\le r}P_i(q)z_1^{(i)}=0$ with $\deg P_i\le d$ exists for $(r,d)\in\{(4,10),
(5,12),(6,14),(8,10)\}$: the coefficient matrix, built from $z_1$ computed to order $4000$, has
full column rank ($55$, $78$, $105$ and $99$ columns respectively) modulo each of two primes.
A relation over $\Q(q)$ may be scaled to a primitive integer one, which survives reduction, so a
full modular column rank does exclude a relation over $\Q$; what it cannot do is reach beyond the
four profiles tested. The honest statement is therefore \emph{not holonomic of order at most $8$
within these degree windows}, and holonomy at higher order or degree remains open. Subject to
that, $z_1$ is an integer power series that is non-algebraic and non-holonomic within the profiles
tested and of radius $<1$, a class for which no arithmetic technology applies: the
Chudnovsky--Andr\'e--Bost theory and the arithmetic holonomy bounds of Calegari--Dimitrov--Tang
\cite{CDT}, which are built for ``integrality plus structure'', all require holonomy, while
P\'olya--Carlson \cite{Carlson1921,Polya1916} requires radius exactly $1$. Integrality alone,
without holonomy, is not currently actionable.

(i) One exact relation ties the lattice across the bases $q$, $-q$, $q^2$: applying the zero
product of Theorem~\ref{thm:zeroproduct-intro} at the three bases and using
$(q;q^2)_\infty(-q;q^2)_\infty=(q^2;q^4)_\infty$,
\[
   \prod_{k\ge1}\varphi_k=1,\qquad
   \varphi_k:=\frac{u_k(q)\,u_k(-q)}{u_k(q^2)} .
\]
The relation is not termwise: $\varphi_1=1+2q^6+4q^8+8q^{10}+18q^{12}+\cdots$,
$\varphi_2=1-2q^6-\cdots$, $\varphi_3=1+2q^{20}+\cdots$, $\varphi_4=1-2q^{28}-\cdots$ (order
$300$; adjacent leading defects cancel in pairs, and the truncated product over $k\le4$ is
$1+O(q^{54})$). No Mahler structure is found. With $u_1$ computed to order $4000$ and every rank
taken modulo each of two primes (a primitive integer relation survives reduction, so full column
rank modulo one prime excludes a rational relation outright), $u_1$ satisfies no polynomial
relation $P\bigl(q,u_1(q),u_1(q^d)\bigr)=0$ for $d=2,3$ at bidegrees $(490,1,1)$, $(150,2,2)$,
$(100,3,3)$, $(60,4,4)$, $(40,5,5)$, $(30,6,6)$, $(24,7,7)$, $(20,8,8)$; no order-two system
relation $P\bigl(q,u_1(q),u_1(q^2),u_1(q^4)\bigr)=0$ at $(30,2,2,2)$, $(20,3,3,3)$; no
difference--differential mix $P\bigl(q,u_1,\theta u_1,u_1(q^2)\bigr)=0$ at $(60,2,2,2)$,
$(40,3,3,3)$, $(30,4,4,4)$; no quasi-modular coupling, in the sense that the six generators
$\{1,\;\theta\!\log u_1(q),\;\theta\!\log
u_1(q^2),\;E_2(q),\;E_2(q^2),\;E_4(q)\}$ admit no $\Q(q)$-linear relation with polynomial
coefficients of degree up to $399$ ($2400$ columns, full rank; so $u_1$ is not a hidden
eta-quotient or $E_2$-ring object of level $2$ within that window); no linear relation among
$\{1,\log u_1(q),\log u_1(q^2)\}$ to degree $665$; and no Mahler relation for the logarithmic
derivative $\theta\!\log u_1$ itself at $(150,2,2)$, $(100,3,3)$, $(60,4,4)$. All twenty-six
profiles are regenerated by \texttt{tools/qcos\_rank}, which prints the row count, the column
count and the rank at each prime. Mahler's method thus finds no foothold on the zero curve at
any complexity reached here; as in (a) and (h), that is a statement about the profiles tested and
not about all profiles.

(j) The fold is the natural point for a rationality reductio, because it is the unique point of
the $(q,z)$-plane carrying \emph{two} vanishings ($G=\theta G=0$; a triple point $G=\theta
G=\theta^2G=0$ is codimension three in two variables, hence absent). The first-jet web there
consists of six values, $A=G(Qz_c)$, $B=\theta G(Qz_c)$, $h=H(z_c)$, $h_1=H(Qz_c)$, $\eta=\theta
H(z_c)$, $\eta_1=\theta H(Qz_c)$, where $H$ is the sine-side companion, which satisfies the mirror
equation $H(z)-a(z)H(Qz)+qH(Q^2z)=0$ (the $z^{1/2}$ in the second solution $z^{1/2}H$ moves the
outer $q$). Exactly two finite-height relations hold: the Casoratian at $z_c$, $Ah=1-q_c$, and its
$\theta$-derivative, $Bh=-A\eta$ (both reproduced to $230$ digits;
\texttt{qcos\_analytic.py}). For any mixed jet-Casoratian
$C_{u,v}(z)=qu(z)v(Qz)-u(Qz)v(z)$ with $L(u)=f$, $\widetilde L(v)=\tilde g$ the $a$-terms cancel
in the increment:
\[
   C_{u,v}(Qz)-C_{u,v}(z)=\tilde g(z)\,u(Qz)-f(z)\,v(Qz),
\]
so every telescoped sum of first-jet orbit bilinears closes back onto a Casoratian. In particular
$B\eta=\sum_{m\ge0}qQ^mz_c\bigl[G\,\theta H-H\,\theta G\bigr](Q^{m+1}z_c)$ (reproduced to
$230$ digits, $246$ terms) is exact but carries an infinite rational series as coefficient: height-infinite.
The ledger of a point evaluation is therefore fixed: each further $\theta$- or
$\delta$-jet level adds at least three new values and exactly one finite-height relation, so the
deficit (four at the fold, five at a generic zero) cannot be closed by jet-raising, and the fold's
two vanishings are the maximum the plane admits. Point evaluation alone, at any point including the
fold, cannot complete a rationality reductio for this module.

(k) All of (c), (f), (g) are instances of a single elementary identity, which closes the
truncation-ledger class outright.

\begin{remark}[denominator--radius inequality; folklore]\label{prop:denrad}
Let $F=\sum_{n\ge0}f_nq^n$ with $f_n\in\Q$, not eventually zero; let $\mathcal D_n$ be the common
denominator of $f_0,\dots,f_n$ and $R$ the radius of convergence. Then
\[
   R\ \le\ \liminf_{n\in S}\mathcal D_n^{1/n},\qquad S:=\{n:f_n\ne0\}.
\]
The restriction to $S$ is necessary. Taking $f_n=c^{-n}$ when $n$ is a power of $2$ and $f_n=0$
otherwise gives $R=c$ while $\liminf_n\mathcal D_n^{1/n}=c^{1/2}$, so the inequality fails over all
$n$ as soon as $c>1$: the long runs of zeros hold $\mathcal D_n$ fixed while $n$ grows. Only the
integral case is used below, and it is unaffected.
Consequently, if $F(s/t)=0$ then the truncation ledger (which requires
$s\,\mathcal D_N^{1/N}<R$) fails for every integer $s\ge1$.
\end{remark}

\noindent This is folklore, not a result of ours. For $\mathcal D_n\equiv1$ it is the standard lemma
that an integer power series of radius $>1$ is a polynomial (Kedlaya \cite[Lem.~11.1.2]{Kedlaya};
the surrounding circle of results is P\'olya--Szeg\H{o} \cite[Part~Eight, Ch.~3]{PolyaSzego}), and
the general case is the same argument after rescaling. Calegari, Dimitrov and Tang
\cite[Lem.~1.1.2]{CDT} name the two halves the \emph{Liouville lower bound} and the \emph{Cauchy
upper bound} and trace the lineage Borel (1894), P\'olya \cite{Polya1916} and Carlson
\cite{Carlson1921}, Dwork \cite{Dwork1960}, Bertrandias, Andr\'e, Bost; the hypothesis that the
denominators grow at most geometrically is Christol's \emph{globally bounded} condition
\cite{Christol1990}, and the $G$-function form of the corollary is in Rivoal
\cite[\S3]{Rivoal2019}. Note also that the inequality is generally \emph{strict}
($f=\sum2^nz^n$ has $\mathcal D_n\equiv1$ and $R=\tfrac12$), so it is not an identity.

\begin{proof}
Whenever $f_n\ne0$ the number $\mathcal D_nf_n$ is a nonzero integer, so $|f_n|\ge1/\mathcal D_n$
and hence
$1/R=\limsup|f_n|^{1/n}\ge\limsup\mathcal D_n^{-1/n}=1/\liminf\mathcal D_n^{1/n}$, which is the
stated bound. The ledger criterion would then force
$\mathcal D_N^{1/N}<R/s\le R\le\liminf\mathcal D_N^{1/N}$, a quantity strictly below its own
$\liminf$. The bound is sharp: $f_n=c^{-n}$ has $\mathcal D_n=c^n$ and $R=c$ exactly.
\end{proof}

The content is that \emph{the denominator floor and the convergence radius are the same quantity}:
buying radius by admitting denominators of growth $c$ costs exactly $c$. Integral series are the
case $\mathcal D_n=1$, recovering (f) and (g); series with geometric denominators gain radius and
lose precisely as much. This is why no rearrangement, reparametrization, or basis change in this
paper's ledgers has ever produced excess, and why none can.

A sharp illustration is available. The radius of $z_1$ is capped by the fold of item (e), where
$z_1$ and $z_2$ collide; but a fold is a square-root branch point of each branch \emph{separately}
and not of the pair, so the elementary symmetric functions of the first $m$ branches (integral
after clearing the poles of $z_k=u_k q^{2(1-k)}$) are analytic through every collision among
them, and the crossing condition survives as the statement that $2q(1-q)$ is a root of
$X^m-e_1X^{m-1}+\cdots$. Their radii increase with $m$. Estimated from coefficient growth at order
$300$, by one and the same estimator for every $m$ so that the rows are comparable, they are
$0.558,\,0.710,\,0.778,\,0.794$ for $m=1,2,3,4$, improving the ledger factor $1/R$ from $1.79$
to $1.26$ (\texttt{qcos\_ledger.py}). The $m=1$ estimate is to be read against the exact value
$|q_c|=0.5546$ from item (e): the estimator reads about $0.5\%$ high there, and the remaining
three rows carry a comparable error, so the improvement is real but the individual figures are
estimates and not certified radii. By Proposition~\ref{prop:denrad} the whole family is capped at
$1$. The $m\to\infty$ limit is the
zero product of Theorem~\ref{thm:zeroproduct-intro}, of radius exactly $1$: the regularized product
is the critical endpoint of this monotone family, which is the precise reason its arithmetic is
confined to the tail.

Together, (c), (f), (g) and Remark~\ref{rem:thetabarrier} express one law from four sides:
for every representation of this problem yet constructed, the arithmetic cost equals the analytic
decay \emph{exactly}, with zero excess; a proof of Conjecture~\ref{conj:Z} must produce strict
excess, and rearrangement, reparametrization, or symmetrization of the object's own
$\Q(q)$-structure provably cannot.
\end{remark}

\subsection{The second-kind forms and the exact reduction}
The $q$-Siegel construction for this class is explicit. For $N\ge1$ put
\[
   P_N(q,z):=(q;q)_{2N}\sum_{k=0}^{N}\frac{(-1)^kq^{k(k-1)}z^k}{(q;q)_{2k}}
   =\sum_{k=0}^{N}(-1)^kq^{k(k-1)}z^k\!\!\prod_{i=2k+1}^{2N}\!\!(1-q^i)\ \in\ \Z[q,z].
\]
\begin{proposition}[the second-kind forms]\label{prop:secondkind}
$\deg_qP_N=2N^2+N$ and $\deg_zP_N=N$ \textup{(exact in $\Z[q,z]$; verified for $N\le8$ by
\texttt{qcos\_ledger.py})}, and the remainder
$R_N(q,z):=(q;q)_{2N}G(q,z)-P_N(q,z)$ satisfies, uniformly for $q$ in compacts of $(0,1)$ and $z$
bounded,
\[
   |R_N(q,z)|=q^{\,N^2(1+o(1))} .
\]
\end{proposition}
\begin{proof}
Membership in $\Z[q,z]$ and $\deg_zP_N=N$ are clear from the displayed product form, since
$(q;q)_{2N}/(q;q)_{2k}=\prod_{i=2k+1}^{2N}(1-q^i)\in\Z[q]$. For the $q$-degree, the $k$-th summand has
degree $k(k-1)+\sum_{i=2k+1}^{2N}i=2N^2+N-k^2-2k$, maximal at $k=0$ with value
$\deg(q;q)_{2N}=2N^2+N$; the maximum is attained once, so no cancellation occurs. For the remainder,
$R_N=(q;q)_{2N}\sum_{k>N}(-1)^kq^{k(k-1)}z^k/(q;q)_{2k}$, whose leading term ($k=N+1$) is
$\pm q^{N(N+1)}z^{N+1}(q;q)_{2N}/(q;q)_{2N+2}$; the tail is dominated by a geometric series in
$q^{2N}$, so $|R_N|=q^{N^2+O(N)}$ uniformly. \textup{(Measured at $q=0.3$, the ratio
$\log|R_N|/(N^2\log q)$ falls to $1.088$ at $N=20$ on the diagonal $z=2q(1-q)$ and to $1.036$ at a
generic $z_0$; \texttt{qcos\_ledger.py}.)}
\end{proof}

\begin{proposition}[reduction to the Siegel margin]\label{prop:reduction}
Let $q=s/t\in(0,1)$ and let $z_0=a/b\in\Q^\times$ with $G(q,z_0)=0$. Then $P_N(q,z_0)=-R_N(q,z_0)$ and
$t^{2N^2+N}b^N\,P_N(q,z_0)\in\Z$. Consequently, if for infinitely many $N$ one has $P_N(q,z_0)\ne0$ and
\[
   |R_N(q,z_0)|<t^{-(2N^2+N)}b^{-N},
\]
then $G(q,z_0)\ne0$: a contradiction. By Proposition~\ref{prop:secondkind} the displayed criterion
reads $q^{N^2}<t^{-2N^2}$ up to $O(N)$ in the exponent, i.e.
\[
   \log(t/s)>2\log t\iff st<1,
\]
which is false. The deficit is exactly the factor $2$ in $\deg_qP_N=\deg(q;q)_{2N}$: the \emph{double}
Pochhammer.
\end{proposition}
\begin{proof}
$P_N(q,z_0)=(q;q)_{2N}G(q,z_0)-R_N(q,z_0)=-R_N(q,z_0)$ by hypothesis; clearing denominators of
$P_N\in\Z[q,z]$ at $q=s/t$, $z=a/b$ gives $t^{\deg_qP_N}b^{\deg_zP_N}P_N(q,z_0)\in\Z$, and a nonzero
integer has modulus $\ge1$. The stated evaluation of the criterion is
Proposition~\ref{prop:secondkind}. Every step of this proof, the clearing included, is
machine-checked (\S\ref{sec:lean}).
\end{proof}

\begin{corollary}[what would suffice]\label{cor:suffices}
Conjecture~\ref{conj:Z} at $z_0$ (hence, on the diagonal, the irrationality of $\beta_2$) follows
from any family of forms $\Lambda_N\in\Z[q,z]$ with $\Lambda_N(q,z_0)\ne0$, remainder
$q^{N^2(1+o(1))}$ against $G$, and
\[
   \deg_q\Lambda_N\;\le\;N^2\,\frac{\log(t/s)}{\log t}\,(1-\varepsilon),
   \qquad\text{in particular }\deg_q\Lambda_N\le N^2(1-\varepsilon)\text{ when }s=1 .
\]
The forms of Proposition~\ref{prop:secondkind} have $\deg_q=2N^2$: it suffices to \emph{halve} the
$q$-degree at the same remainder, or equivalently to double the remainder decay to $q^{2N^2}$. A
single-Pochhammer clearing $(q;q)_N$ would have degree $\tfrac12N^2(1+o(1))$ and would suffice at once;
no such representation of $G$ along a non-$q$-power argument exists in the
Rogers--Ramanujan/Bailey framework.
\end{corollary}

\begin{remark}[the theta-decay barrier: why Corollary~\ref{cor:suffices} is unreachable by truncation]
\label{rem:thetabarrier}
The hypothesis of Corollary~\ref{cor:suffices} cannot be met by \emph{any} truncate-and-clear form.
Let $\sum_kc_k(q)q^{\alpha k^2}z^k$ ($\alpha>0$, $c_k\in\Q(q)$) be truncated at $N$ and cleared. The
$k=N$ summand already contributes $q^{\alpha N^2}$, so $\deg_q\ge\alpha N^2$; the remainder is the
first omitted term, of size $q^{\alpha N^2(1+o(1))}$. The criterion $|R_N|<t^{-\deg_q}$ therefore reads
\[
   \alpha N^2\log(t/s)\;>\;\alpha N^2\log t\iff s<1,
\]
false for every $q=s/t\in(0,1)$: \emph{the decay and the degree are driven by the same factor
$q^{\alpha k^2}$ and cancel exactly}. Denominators only make matters worse. The cleared truncations
have $\deg_q=N^2$ exactly for the theta series $\sum q^{k^2}z^k$ and for the Rogers--Ramanujan
series $\sum q^{k^2}/(q;q)_k$, so both sit at $\beta:=\deg_q/N^2=1$, the boundary; the Hahn--Exton
series sits at $\beta=2+1/N$, that is at $2.25$, $2.17$, $2.12$ for $N=4,6,8$, a further factor $2$
past a boundary that is itself already unreachable (\texttt{qcos\_ledger.py}).

This describes the shape of the known landscape and not merely our position in it. The
transcendence of theta values and of Rogers--Ramanujan values, both at $\beta=1$, is known
\emph{only} through Nesterenko's modular method \cite{Nesterenko} and never through a
Siegel/Liouville margin, precisely because no theta-decay series can cross $\beta<1$. The two
mechanisms left for Conjecture~\ref{conj:Z} are therefore modularity and a multiplicity theory for
the associated $(\sigma,\delta)$-module. For the first, no analogue of \cite{Nesterenko} is known
off the modular class, and we have not found $G$ inside it; that is the state of knowledge, not a
proof that $G$ lies outside (see \cite{Bonfioli2}). The double Pochhammer is not the whole
obstruction: it is a factor $2$ on top of a structural one.
\end{remark}

\subsection{An unconditional effective bound}
Although Conjecture~\ref{conj:Z} is open, the specific root $q^\ast$ admits an effective
non-rationality statement, proved by certified interval arithmetic.

\begin{proposition}[effective non-rationality of $q^\ast$]
\label{prop:effective}
The least positive root $q^\ast$ of $S$ admits no rational representation $s/t$ with
$t\le10^{1044}$. Consequently, if $\beta_2=u/v$ in lowest terms then $u>10^{522}$. The method is
not new: it is Brent's \cite{Brent1977}, who proved by the same continued-fraction argument that a
rational representation of Euler's constant needs a large denominator, and it rests on Khinchin's
Theorem~17 (every convergent is a best approximation of the second kind)
\cite[Thm.~17]{Khinchin}. Lagarias \cite[\S3.15]{Lagarias2013} surveys the Brent--McMillan
denominator bounds for Euler's constant and quotes Khinchin's theorem in their proof. What is new
here is only the bound for this particular constant.
\end{proposition}

\begin{proof}
All steps are certified interval arithmetic (outward rounding; \texttt{mpmath.iv}; script
\texttt{beta2\_effective\_irrationality.py}, which bootstraps $q^\ast$ from a $25$-digit start by
Newton at doubling precision and so needs no stored input), with the series tail bounded by
$(q;q)_{2k}\ge(1-2q)/(1-q)>0$ for $q<\tfrac12$ and term ratios $<\tfrac12$ beyond the truncation.
\textup{(i)} $S>0$ on $[0,0.40]$ ($400$ certified chunks). \textup{(ii)} $S'\le-1$ on $[0.40,0.46]$
($600$ certified chunks; $S'$ evaluated term-by-term with the logarithmic-derivative weights), so
$S$ has at most one zero in $[0.40,0.46]$. \textup{(iii)} with $c$ the
$2090$-digit truncation of $q^\ast$, $a=c-2\cdot10^{-2090}$ and $b=c+2\cdot10^{-2090}$, certified
evaluation gives $S(a)>0>S(b)$ (working precision $2130$ digits, series tail enclosed in
$\pm10^{-2125}$); by \textup{(i)}--\textup{(ii)}, $q^\ast\in[a,b]$. \textup{(iv)} the interval
$[a,b]$ contains the rational $s_0/t_0$ (Stern--Brocot construction) with
$9.5\cdot10^{1044}<t_0<10^{1045}$, verified by exact rational comparison, and certified
$S(s_0/t_0)\ne0$. If $m/n\in[a,b]$ with $n\le10^{1044}$ and $m/n\ne s_0/t_0$, then
$|m/n-s_0/t_0|\ge1/(nt_0)>10^{-2089}>b-a=4\cdot10^{-2090}$, impossible; so $m/n=s_0/t_0$, whence
$n=t_0>10^{1044}$, a contradiction. Thus no rational with denominator $\le10^{1044}$ lies in
$[a,b]$ at all, in particular none equals $q^\ast$. For $\beta_2=u/v$ in lowest terms,
$q^\ast=v^2/u^2$ in lowest terms, so $u^2>10^{1044}$. The Diophantine half of this argument,
steps \textup{(iv)} onward, is machine-checked, with steps \textup{(i)}--\textup{(iii)} as
hypotheses (\S\ref{sec:lean}).
\end{proof}


\section{Atom decomposition and status}\label{sec:atoms}

The table below is the paper's status register. Every numbered statement, and every remark
carrying mathematical content, appears exactly once, with a one-line paraphrase, its label, the
statements it depends on, and its status. \textsc{proved} means a complete human-style proof is
written out here or cited from the literature; \textsc{heuristic} means the support is
computational; \textsc{conjecture} means believed and unproved. A statement is exactly as strong
as the weakest label in its dependency chain, and the table is arranged so that this can be
checked by reading upwards. A star marks a statement at \textsc{proved} whose formalization is
incomplete; \S\ref{sec:lean} records, for each starred entry, the coverage attained and what
blocks the rest. There are no unlabelled mathematical statements in the paper and no atom without
a path to Conjecture~\ref{conj:Z} or to Proposition~\ref{prop:effective}.

\begin{center}
\scriptsize
\begin{tabular}{@{}l>{\raggedright\arraybackslash}p{5.4cm}>{\raggedright\arraybackslash}p{2.4cm}>{\raggedright\arraybackslash}p{1.7cm}>{\raggedright\arraybackslash}p{2.2cm}@{}}
\hline
ID & statement & label & depends on & status\\
\hline
A1 & Thm~\ref{thm:main-transc}: $G(q,\cdot)$ satisfies no algebraic differential equation over
  $\C(z)$ & \textsc{proved} & \cite{ADH2021} & done$^\star$, known\\
A2 & Lem~\ref{lem:irreducible}: the rank-$2$ $\sigma$-module is irreducible; the Riccati equation
  has no rational solution & \textsc{proved} & none & done$^\star$\\
A3 & Thm~\ref{thm:sl2}: the difference Galois group contains $\mathrm{SL}_2$ & \textsc{proved} &
  A2, \cite{RSZ} & done$^\star$\\
A4 & Prop~\ref{prop:nonintegrable}: no $\sigma\delta$-integrability relation, not even with a
  scalar twist & \textsc{proved} & none & done$^\star$\\
A5 & Cor~\ref{cor:hypertranscendence}: $G(q,\cdot)$ is differentially transcendental, by the
  parametrized route & \textsc{proved} & A3, A4, \cite{DHR2018} & done$^\star$\\
A6 & Thm~\ref{thm:zeroproduct} (= Thm~\ref{thm:zeroproduct-intro}): $C_1=1/(q;q)_\infty$ and
  $\prod_kz_kQ^{k-1}=(q;q^2)_\infty$ & \textsc{proved} & none & done$^\star$\\
A7 & Rem~\ref{rem:zeroproduct-dioph}: if $z^\ast$ is algebraic then the tail product is
  transcendental & \textsc{proved} & A6, \cite{Nesterenko} & done$^\star$\\
A8 & Thm~\ref{thm:integrality}: the Hensel root $u_k$ of $F_k$ lies in $1+q\Z[[q]]$ &
  \textsc{proved} & A6 & done$^\star$; the identification of $u_k$ with $q^{2(k-1)}z_k$ needs
  convergence, hence A22\\
A9 & Lem~\ref{lem:swap}: the parameter--argument swap for ${}_1\phi_1(0;b;Q,z)$ &
  \textsc{proved} & \cite{KS1992} & done$^\star$, known\\
A10 & Prop~\ref{prop:latticevalues}: $\ord_qG(q,Q^{-m})=(m+1)^2$ with leading coefficient
  $(-1)^{m+1}$ & \textsc{proved} & A9, \cite{KS1992} & done$^\star$, known\\
A11 & Cor~\ref{cor:hexagonal}: $\operatorname{onset}(u_k-1)=k(2k-1)$, leading coefficient $-1$ &
  \textsc{proved} & A8, A10 & done$^\star$\\
A12 & Rem~\ref{rem:latticevalues}: order reflection $-\tfrac12\mapsto m+1$; backward-orbit
  constant; vanishing transferred to the parameter slot & \textsc{proved} & A9, A10 &
  done$^\star$\\
A13 & Prop~\ref{prop:sinelattice}: the sine-side integrality, lattice values and onsets
  $k(2k+1)$ & \textsc{proved} & A8, A9 & done$^\star$\\
A14 & Rem~\ref{rem:triangular}: the two onset families exhaust the triangular numbers, each once
  & \textsc{proved} & A11, A13 & done$^\star$\\
A15 & Prop~\ref{prop:stablelaw}: both deviation families stabilize $q$-adically to
  $1/(q^2;q^2)_\infty^2$ & \textsc{proved} & A10, A11, A13 & done$^\star$\\
A16 & Rem~\ref{prop:denrad}: $R\le\liminf_{n\in S}\mathcal D_n^{1/n}$, and the truncation ledger
  fails for every integer $s\ge1$ & \textsc{proved} & none & done$^\star$, folklore\\
A17 & Prop~\ref{prop:secondkind}: $\deg_qP_N=2N^2+N$, $\deg_zP_N=N$,
  $|R_N|=q^{N^2(1+o(1))}$ & \textsc{proved} & none & done$^\star$\\
A18 & Prop~\ref{prop:reduction}: the reduction to the Siegel margin, and its failure by the
  factor $2$ & \textsc{proved} & A17 & done$^\star$\\
A19 & Cor~\ref{cor:suffices}: halving the $q$-degree at the same remainder would suffice &
  \textsc{proved} & A18 & done$^\star$\\
A20 & Rem~\ref{rem:thetabarrier}: no truncate-and-clear form meets A19's hypothesis &
  \textsc{proved} & A19 & done$^\star$\\
A21 & Rem~\ref{rem:zerocurve-arith}, items (c), (f), (g), (j), (k): the linear-height ledger, the
  uniformization cap, the discreteness cap, the jet deficit at the fold, and their common source
  & \textsc{proved} & A8, A16 & done$^\star$\\
A22 & Rem~\ref{rem:zerocurve-arith}, items (b), (d), (e): the radius equals $|q_c|$; the fold
  coordinates and the fold constant; the coefficient sequence is new &
  \textsc{heuristic} & A8 & numerical, $230$ digits\\
A23 & Rem~\ref{rem:zerocurve-arith}, items (a), (h), (i): the exclusions of algebraic, holonomic,
  Mahler, system, quasi-modular and log-linear relations &
  \textsc{proved} within the profiles listed, \textsc{heuristic} beyond them & A8 &
  rank certificates\\
A24 & Prop~\ref{prop:effective}: $q^\ast$ has no rational representation of denominator
  $\le10^{1044}$ & \textsc{proved} & A25 & done$^\star$\\
A25 & The certified enclosure of $q^\ast$ and the Stern--Brocot fraction $s_0/t_0$ &
  \textsc{proved} & none & interval arithmetic; hypothesis of A24\\
A26 & Conj~\ref{conj:Z}: every positive zero $z_k(q)$ is transcendental at algebraic $q$ &
  \textsc{conjecture} & none & open; the goal\\
\hline
\end{tabular}
\end{center}

\noindent Reading of the table. A26 is the goal; A1--A5 place the function outside the reach of
differential-algebraic methods, A6--A15 supply the exact structure of the zero lattice, and
A16--A24 map the negative space: which mechanisms cannot reach A26 and what the best
unconditional statement about the distinguished zero currently is. A22 and A23 are the only
entries not at \textsc{proved}, and neither is used in the proof of any other entry. A8 is
stated for the formal Hensel root and is unconditional as such; only its identification with the
analytic $k$-th zero passes through A22, and that qualification is carried in the theorem itself
rather than in a remark. A23 is terminal. A25 is separated from A24 because it is exactly the part of
Proposition~\ref{prop:effective} that a machine check consumes as a hypothesis rather than proves;
\S\ref{sec:lean} says so as well.

\section{Machine verification}\label{sec:lean}

The arithmetic content of this paper is formalised in Lean~4 (toolchain \texttt{v4.30.0}, Mathlib
\texttt{v4.30.0}). Five files in \texttt{lean/with\_mathlib/} of the cited repository carry it:
\begin{center}
\texttt{QCosineExponents}, \texttt{QCosineLattice}, \texttt{QSiegelLedger},
\texttt{QEffectiveExclusion}, \texttt{QZeroSeries}.
\end{center}
\noindent All five are registered in
\texttt{lakefile.toml} and build with a bare \texttt{lake build}. Every declaration named below is
\texttt{sorry}-free, and its axiom list was read off a cold build with \texttt{\#print axioms}. The
only axioms occurring are Lean's standard \texttt{propext}, \texttt{Classical.choice} and
\texttt{Quot.sound}, with the single exception of \texttt{QZeroSeries.lean}, where each theorem is
decided by compiled evaluation and therefore also carries the reflection axiom that
\texttt{native\_decide} generates.

\begin{center}
\scriptsize
\begin{tabular}{@{}ll>{\raggedright\arraybackslash}p{7.2cm}@{}}
\hline
Statement & Coverage & Lean declarations\\
\hline
Thm~\ref{thm:main-transc} & none & (see below)\\
Thm~\ref{thm:zeroproduct-intro} & exponent steps & as Thm~\ref{thm:zeroproduct}\\
Lem~\ref{lem:irreducible} & none & (see below)\\
Thm~\ref{thm:sl2} & exponent gap & \texttt{supermult\_exponent}, \texttt{gap\_ge},
  \texttt{gap\_eq\_at\_ends}\\
Prop~\ref{prop:nonintegrable} & none & (see below)\\
Cor~\ref{cor:hypertranscendence} & none & (see below)\\
Thm~\ref{thm:zeroproduct} & exponent steps & \texttt{zeroproduct\_reindex},
  \texttt{zeroproduct\_peak\_bound}, \texttt{zeroproduct\_peak\_slack}, \texttt{euler\_split}\\
Thm~\ref{thm:integrality} & inputs, and order $q^{80}$ & \texttt{integrality\_exponent},
  \texttt{newtonExp\_nonneg}, \texttt{newtonExp\_eq\_zero\_iff},
  \texttt{newtonExp\_second\_difference}, \texttt{newton\_residual\_shape},
  \texttt{residualPoly\_eval\_one}, \texttt{residualPoly\_derivative\_eval\_one},
  \texttt{residualPoly\_derivative\_isUnit}, \texttt{qPochhammer\_isUnit},
  \texttt{u1\_coefficients}, \texttt{u2\_prefix}, \texttt{cosine\_residual\_vanishes}\\
Lem~\ref{lem:swap} & exponent identities & \texttt{swap\_symmetry\_exponent},
  \texttt{swap\_base\_collapse}\\
Prop~\ref{prop:latticevalues} & no-cancellation step & \texttt{lattice\_order\_gap},
  \texttt{lattice\_order\_strictMono}, \texttt{lattice\_order\_leading}\\
Cor~\ref{cor:hexagonal} & onsets, $k\le6$ & \texttt{hexagonal\_onset},
  \texttt{hexagonal\_next\_correction}, \texttt{cosine\_onset}, \texttt{cosine\_onsets}\\
Prop~\ref{prop:sinelattice} & onsets, $k\le5$ & \texttt{sine\_shift}, \texttt{sine\_onset},
  \texttt{onsets\_differ}, \texttt{onsets\_consecutive}, \texttt{sine\_onsets},
  \texttt{sine\_residual\_vanishes}\\
Prop~\ref{prop:stablelaw} & depths, $k\le6$ & \texttt{stablelaw\_numerator\_depth},
  \texttt{stablelaw\_numerator\_depth\_sharp}, \texttt{stablelaw\_pair\_exponent},
  \texttt{stablelaw\_pair\_weight}, \texttt{stablelaw\_pair\_parity},
  \texttt{stablelaw\_denominator\_depth}, \texttt{twoColoured\_prefix},
  \texttt{cosine\_agreement\_depth}, \texttt{sine\_agreement\_depth}\\
Rem~\ref{prop:denrad} & Liouville half & \texttt{one\_le\_abs\_of\_ne\_zero}\\
Prop~\ref{prop:secondkind} & degrees & \texttt{tail\_sum}, \texttt{degree\_ledger},
  \texttt{degree\_at\_zero}, \texttt{degree\_lt\_of\_pos}\\
Prop~\ref{prop:reduction} & complete & \texttt{prod\_one\_sub\_div}, \texttt{clearedTerm\_eq},
  \texttt{cleared\_eq}, \texttt{reduction}, \texttt{margin\_iff}, \texttt{margin\_fails},
  \texttt{st\_ge\_one}\\
Cor~\ref{cor:suffices} & degree threshold at the coefficient $2$ & \texttt{margin\_iff}
  (see below)\\
Rem~\ref{rem:thetabarrier} & complete & \texttt{theta\_barrier}\\
Prop~\ref{prop:effective} & complete & \texttt{farey\_gap},
  \texttt{unique\_rational\_in\_interval}, \texttt{no\_rational\_of\_denominator\_le},
  \texttt{denominator\_gt\_of\_mem}, \texttt{width\_lt\_gap},
  \texttt{qstar\_no\_small\_denominator}\\
\hline
\end{tabular}
\end{center}

\medskip
Two entries deserve comment because the formalisation changed the text. The displayed derivation
of Proposition~\ref{prop:sinelattice}(c) previously omitted the prefactor $q^{k(k-1)}$ and so
evaluated to $1$ rather than to $k(2k+1)$; and the numerator step of
Proposition~\ref{prop:stablelaw} previously gave the relative order of the $r$-th term as
$(k+r)^2-k(2k-1)$, which at $k=2$, $r=1$ is $3$ and contradicts the depth $2k+1=5$ asserted two
lines later. Both are corrected above; neither conclusion changes.

\medskip
\noindent The following are \emph{not} formalised.

\begin{itemize}
\item \emph{All of \S\ref{sec:galois}.} Theorem~\ref{thm:sl2},
Proposition~\ref{prop:nonintegrable}, Corollary~\ref{cor:hypertranscendence} and
Theorem~\ref{thm:main-transc} are statements of difference Galois theory and of its parametrized
refinement. Mathlib contains no difference Galois theory, no $q$-difference modules and no
Picard--Vessiot theory for difference equations, so there is nothing to reduce them to;
formalising them means formalising \cite{HS} and \cite{DHR2018} first. Only the exponent gap that
replaced the numerical check in the proof of Theorem~\ref{thm:sl2} is machine-checked. A file
whose content was ``assume \cite[Prop.~2.9]{DHR2018}, conclude'' would add nothing and is
deliberately not supplied. The same applies to Lemma~\ref{lem:irreducible}, whose argument is a
valuation-theoretic exhaustion over the $Q$-orbits of $\C^\ast$.
\item \emph{The analytic steps of Theorem~\ref{thm:zeroproduct}}: the uniform peak estimate, the
Rouch\'e localisation and the dominated-convergence passage. The exponent identities those steps
consume are formalised; the estimates are not. Mathlib has the dominated convergence theorem but
no Rouch\'e theorem, so step~(3) has no target to reduce to at present.
\item \emph{The Hensel step of Theorem~\ref{thm:integrality}.} Mathlib does supply
\texttt{IsAdicComplete (Ideal.span \{X\}) (PowerSeries R)} together with
\texttt{IsAdicComplete.henselianRing}, so $\Z[[q]]$ is Henselian at $(q)$ in the library's sense.
That instance does not apply here: \texttt{HenselianRing.is\_henselian} requires a \emph{monic
polynomial}, whereas $F_k$ is a power series in $u$ whose top behaviour is not a unit. The missing
ingredient is a Weierstrass-preparation form of the implicit function theorem for $A[[u]]$ over an
adically complete $A$ that is not local, and Mathlib's Weierstrass preparation assumes $A$ local.
What is verified instead is the finite-order content, which is the part the paper displays: the
Newton recursion carried out entirely over $\Z$, with the residual $F_k(q,u_k)$ checked to vanish
to order $q^{80}$ for $k\le6$ on the cosine side and $k\le5$ on the sine side.
\item \emph{The Cauchy half of Remark~\ref{prop:denrad}}, namely
$1/R=\limsup|f_n|^{1/n}$. The Liouville half, $|f_n|\ge1/\mathcal D_n$ on the support, is
formalised; the radius identity is Cauchy--Hadamard and is used as such.
\item \emph{The remainder estimate of Proposition~\ref{prop:secondkind}}, $|R_N|=q^{N^2(1+o(1))}$.
It is an estimate on a $q$-series tail and enters \texttt{reduction} as a hypothesis, which is the
honest place for it: what is proved is that the estimate, together with $P_N(q,z_0)\ne0$, forces a
contradiction.
\item \emph{The interval arithmetic of Proposition~\ref{prop:effective}}: the enclosure of
$q^\ast$, the Stern--Brocot rational $s_0/t_0$ and the certified sign changes come from
\texttt{beta2\_effective\_irrationality.py} and enter \texttt{qstar\_no\_small\_denominator} as
hypotheses. The Diophantine argument that consumes them, including the width comparison
$4\cdot10^{-2090}<1/(10^{1044}t_0)$ for every $t_0<10^{1045}$, is formalised in full.
\item \emph{The general degree threshold of Corollary~\ref{cor:suffices}.} Its content is the
comparison already formalised as \texttt{margin\_iff}, which decides
$2\log t<\log(t/s)\iff st<1$: that is the threshold at the coefficient $2$ realised by
Proposition~\ref{prop:secondkind}. What is not formalised is the same comparison with the
coefficient $2$ replaced by the parameter $\deg_q\Lambda_N/N^2$, together with the $o(1)$
bookkeeping in the remainder hypothesis. No Mathlib gap blocks this: the threshold is one
inequality between logarithms and the asymptotic hypothesis is expressible with
\texttt{Filter.Tendsto}. It is open formalization work, and it is recorded here as such rather
than left without a status.
\item \emph{The exclusion searches of Remark~\ref{rem:zerocurve-arith}}, which are rank
computations on large integer matrices over $\Z/p$. Each is a certified rank witness and nothing
in Mathlib blocks formalising it; what would be needed is to import the coefficient data and the
elimination transcript, which is a data-volume question and not a mathematical one. They are
reported as computations, regenerated by \texttt{tools/qcos\_rank}.
\end{itemize}

\medskip
\noindent One point of method is worth recording, since it is what made
Proposition~\ref{prop:reduction} cheap to verify. Denominator clearing is normally stated for a
general two-variable integer polynomial, in terms of its two total degrees, and proving it in that
form is substantial. It is not needed. Since $P_N$ is given by an explicit product formula, the
cleared quantity can be written down as an integer,
\[
   t^{2N^2+N}b^N P_N(s/t,a/b)\;=\;\sum_{k=0}^{N}(-1)^ks^{k(k-1)}a^kb^{N-k}t^{k^2+2k}
   \!\!\prod_{i=2k+1}^{2N}\!\!\bigl(t^i-s^i\bigr),
\]
and the whole content becomes the exponent ledger
$k(k-1)+\sum_{i=2k+1}^{2N}i+(k^2+2k)=2N^2+N$, which is \texttt{degree\_ledger}.

\section*{Declarations}
\noindent\textbf{Funding.} No funding was received for this work.
\smallskip\\
\noindent\textbf{Competing interests.} The author declares no competing interests.
\smallskip\\
\noindent\textbf{Data availability.} All code, numerical certification data and Lean sources
supporting the computational and machine-checked statements are openly available in the repository
cited on the title page and archived on Zenodo (concept DOI
\texttt{10.5281/zenodo.20370090}). The Lean files are listed in \S\ref{sec:lean}.
\smallskip\\
\noindent\textbf{Use of AI tools.} Large language models (Anthropic Claude) were used in the
preparation of this manuscript, both in drafting and in computation. Every computational claim
is certified by the scripts in the cited repository, independently of any model output. The
author takes full responsibility for the content.

\begin{thebibliography}{99}
\bibitem{AKV2002} M.~Amou, M.~Katsurada, K.~V\"a\"an\"anen, \emph{On the values of certain
$q$-hypergeometric series II}, in: Analytic Number Theory, eds.\ C.~Jia, K.~Matsumoto,
Developments in Mathematics~\textbf{6}, Kluwer Acad.\ Publ., Boston, MA, 2002, 17--25.
\bibitem{Brent1977} R.~P.~Brent, \emph{Computation of the regular continued fraction for Euler's
constant}, Math.\ Comp.\ \textbf{31} (1977), 771--777.
\bibitem{CDT} F.~Calegari, V.~Dimitrov, Y.~Tang, \emph{The unbounded denominators conjecture},
J.~Amer.\ Math.\ Soc.\ \textbf{38} (2025), no.~3, 627--702; arXiv:2109.09040.
\bibitem{Kedlaya} K.~S.~Kedlaya, \emph{Weil cohomology in practice}, lecture notes,
\url{https://kskedlaya.org/weil-cohom/}, Ch.~11 (Dwork's proof of rationality), Lemma~11.1.2.
These are living notes with no version stamp; the statement used is that an integer power series
with radius of convergence $>1$ is a polynomial, whose two-line proof is reproduced above.
\bibitem{Lagarias2013} J.~C.~Lagarias, \emph{Euler's constant: Euler's work and modern developments},
Bull.\ Amer.\ Math.\ Soc.\ \textbf{50} (2013), 527--628.
\bibitem{LorchMuldoon} L.~Lorch, M.~E.~Muldoon, \emph{Transcendentality of zeros of higher
derivatives of functions involving Bessel functions}, Int.\ J.\ Math.\ Math.\ Sci.\ \textbf{18}
(1995), 551--560.
\bibitem{Rivoal2019} T.~Rivoal, \emph{Les $E$-fonctions et $G$-fonctions de Siegel}, in:
Journ\'ees math\'ematiques X-UPS, Ann\'ee 2019: P\'eriodes et transcendance, \'Ed.\ P.~Harinck,
A.~Plagne, C.~Sabbah, Les \'Editions de l'\'Ecole polytechnique / Centre Mersenne, 2019, 197--298.
\bibitem{Lindemann1882} F.~Lindemann, \emph{Ueber die Zahl $\pi$},
Math.\ Ann.\ \textbf{20} (1882), 213--225.

\bibitem{Siegel1929} C.~L.~Siegel, \emph{\"Uber einige Anwendungen diophantischer Approximationen},
Abh.\ Preuss.\ Akad.\ Wiss., Phys.-math.\ Kl.\ (1929), Nr.~1, 70 pp.; Gesammelte Abhandlungen~I,
Springer, 1966, 209--266. English translation by C.~Fuchs, ed.\ U.~Zannier, Quaderni della Scuola
Normale Superiore di Pisa, Monographs~2, Edizioni della Normale, 2014, 81--138.
\bibitem{ADH2021} B.~Adamczewski, T.~Dreyfus, C.~Hardouin, \emph{Hypertranscendence and linear
difference equations}, J.~Amer.\ Math.\ Soc.\ \textbf{34} (2021), 475--503.
\bibitem{ABC2003} L.~D.~Abreu, J.~Bustoz, J.~L.~Cardoso, \emph{The roots of the third Jackson
$q$-Bessel function}, Int.\ J.\ Math.\ Math.\ Sci.\ \textbf{2003}, no.~67, 4241--4248.
\bibitem{AMV} M.~Amou, T.~Matala-aho, K.~V\"a\"an\"anen, \emph{On Siegel--Shidlovskii's theory for
$q$-difference equations}, Acta Arith.\ \textbf{127} (2007), 309--335.
\bibitem{ADR} C.~E.~Arreche, T.~Dreyfus, J.~Roques, \emph{Differential transcendence criteria for
second-order linear difference equations and elliptic hypergeometric functions}, J.~\'Ec.\ polytech.\
Math.\ \textbf{8} (2021), 147--168.
\bibitem{Bonfioli2} V.~Bonfioli, \emph{Transcendence of a planar reflection--group growth series and
its relaxed companion}, preprint (2026), archived at \texttt{doi:10.5281/zenodo.20370090}
(file \texttt{paper/journal/paper2.pdf}).
\bibitem{DHR2018} T.~Dreyfus, C.~Hardouin, J.~Roques, \emph{Hypertranscendence of solutions of Mahler
equations}, J.~Eur.\ Math.\ Soc.\ \textbf{20} (2018), 2209--2238.
\bibitem{HS} C.~Hardouin, M.~F.~Singer, \emph{Differential Galois theory of linear difference
equations}, Math.\ Ann.\ \textbf{342} (2008), 333--377.
\bibitem{IsmailZhang2007} M.~E.~H.~Ismail, C.~Zhang, \emph{Zeros of entire functions and a problem
of Ramanujan}, Adv.\ Math.\ \textbf{209} (2007), 363--380.
\bibitem{KS1992} T.~H.~Koornwinder, R.~F.~Swarttouw, \emph{On $q$-analogues of the Fourier and Hankel
transforms}, Trans.\ Amer.\ Math.\ Soc.\ \textbf{333} (1992), 445--461.
\bibitem{Morita2011} T.~Morita, \emph{A connection formula of the Hahn--Exton $q$-Bessel function},
SIGMA \textbf{7} (2011), 115, 11 pp.
\bibitem{Nishioka2016} S.~Nishioka, \emph{Proof of unsolvability of $q$-Bessel equation using
valuations}, J.~Math.\ Sci.\ Univ.\ Tokyo \textbf{23} (2016), 763--789.
\bibitem{KS} H.~T.~Koelink, R.~F.~Swarttouw, \emph{On the zeros of the Hahn--Exton $q$-Bessel function
and associated $q$-Lommel polynomials}, J.\ Math.\ Anal.\ Appl.\ \textbf{186} (1994), 690--710.
\bibitem{Nesterenko} Yu.~V.~Nesterenko, \emph{Modular functions and transcendence questions},
Sb.\ Math.\ \textbf{187} (1996), 1319--1348.
\bibitem{Carlson1921} F.~Carlson, \emph{\"Uber Potenzreihen mit ganzzahligen Koeffizienten},
Math.\ Z.\ \textbf{9} (1921), 1--13.
\bibitem{Christol1990} G.~Christol, \emph{Globally bounded solutions of differential equations}, in:
Analytic Number Theory (Tokyo, 1988), Lecture Notes in Math.\ \textbf{1434}, Springer, Berlin, 1990,
45--64.
\bibitem{Dwork1960} B.~Dwork, \emph{On the rationality of the zeta function of an algebraic variety},
Amer.\ J.\ Math.\ \textbf{82} (1960), 631--648.
\bibitem{Exton1978} H.~Exton, \emph{A basic analogue of the Bessel--Clifford equation},
J\~n\=an\=abha \textbf{8} (1978), 49--56.
\bibitem{GR} G.~Gasper, M.~Rahman, \emph{Basic Hypergeometric Series}, 2nd ed., Encyclopedia of
Mathematics and its Applications~\textbf{96}, Cambridge University Press, Cambridge, 2004.
\bibitem{Hahn1949} W.~Hahn, \emph{Beitr\"age zur Theorie der Heineschen Reihen. Die 24 Integrale der
hypergeometrischen $q$-Differenzengleichung. Das $q$-Analogon der Laplace-Transformation},
Math.\ Nachr.\ \textbf{2} (1949), no.~6, 340--379.
\bibitem{Hayman1956} W.~K.~Hayman, \emph{A generalisation of Stirling's formula},
J.~Reine Angew.\ Math.\ \textbf{196} (1956), 67--95.
\bibitem{Jacobi1829} C.~G.~J.~Jacobi, \emph{Fundamenta Nova Theoriae Functionum Ellipticarum},
Sumtibus fratrum Borntr\"ager, Regiomonti, 1829; Gesammelte Werke~I, Reimer, Berlin, 1881.
\bibitem{Khinchin} A.~Ya.~Khinchin, \emph{Continued Fractions}, transl.\ Scripta Technica,
ed.\ H.~Eagle, University of Chicago Press, Chicago, 1964; reprinted Dover, Mineola NY, 1997.
\bibitem{Polya1916} G.~P\'olya, \emph{\"Uber Potenzreihen mit ganzzahligen Koeffizienten},
Math.\ Ann.\ \textbf{77} (1916), 497--513.
\bibitem{PolyaSzego} G.~P\'olya, G.~Szeg\H{o}, \emph{Problems and Theorems in Analysis~II},
transl.\ C.~E.~Billigheimer, Springer, Berlin, 1976; reprinted in Classics in Mathematics, 1998.
\bibitem{RSZ} J.-P.~Ramis, J.~Sauloy, C.~Zhang, \emph{Local analytic classification of
$q$-difference equations}, Ast\'erisque \textbf{355}, Soc.\ Math.\ France, 2013.
\bibitem{vdPS} M.~van der Put, M.~F.~Singer, \emph{Galois Theory of Difference Equations},
Lecture Notes in Mathematics~\textbf{1666}, Springer, Berlin, 1997.
\end{thebibliography}

\end{document}
